\exerciseshdr{Exercises}

\exercisegroup

\begin{enumerate}
\def\labelenumi{\arabic{enumi})}
\tightlist
\item
  Let \((\mathrm{Q}_{i})_{i\in\mathrm{I}}\) be a family of quivers. We denote by S the sum set of the family \((\mathrm{Som}(\mathrm{Q}_{i}))\), by F the sum set of the family \((\mathrm{Fl}(\mathrm{Q}_{i}))\), and by o and t the maps from F to S induced by the origin and target maps of the quivers \(Q_{i}\).
\end{enumerate}

\begin{enumerate}
\def\labelenumi{\alph{enumi})}
\item
  Prove that \(\mathrm{Q} = (\mathrm{S}, \mathrm{F}, o, t)\) is a quiver. We call it the sum quiver of the family \((\mathrm{Q}_{i})\).
\item
  For \(i \in I\) there exists a unique quiver morphism \(j_{i}\) from \(Q_{i}\) to Q such that the map \(\operatorname{Som}(j_{i})\) is the canonical map from \(\operatorname{Som}(Q_{i})\) into S and the map \(\operatorname{Fl}(j_{i})\) is the canonical map from \(\operatorname{Fl}(Q_{i})\) into S.
\item
  Let K be a quiver and, for every \(i \in I\), let \(\varphi_{i}\) be a quiver morphism from \(Q_{i}\) to K. Show that there exists a unique quiver morphism \(\varphi\) from Q to K such that one has \(\varphi \circ j_{i} = \varphi_{i}\) for all \(i \in I\).
\end{enumerate}

\begin{enumerate}
\def\labelenumi{\arabic{enumi})}
\setcounter{enumi}{1}
\tightlist
\item
  Let \(Q = (S, F, o, t)\) be a quiver. For every \(s \in S\), we call \(d_{+}(s)\) (resp. \(d_{-}(s)\)) the outgoing (resp. incoming) degree of Q at s and define \(d_{+}(s)\) (resp. \(d_{-}(s)\)) as the cardinality of the set \(\{f \in F \mid o(f) = s\}\) (resp. \(\{f \in F \mid t(f) = s\}\)). We say that Q is locally finite if \(d_{+}(s)\) and \(d_{-}(s)\) are finite for every \(s \in S\), and that it is finite if the sets F and S are finite.
\end{enumerate}

Prove the equalities. \(\sum_{s\in S}d_{+}(s) = \sum_{s\in S}d_{-}(s) = \mathrm{Card}(\mathrm{F})\)

\begin{enumerate}
\def\labelenumi{\arabic{enumi})}
\setcounter{enumi}{2}
\tightlist
\item
  Let \(\mathrm{Q} = (\mathrm{S},\mathrm{F},o,t)\) be a locally finite quiver. We define the matrix \(M_{\mathrm{Q}} = (m_{i,j})\) of type (S,S) with elements in \(\mathbf{Z}\), by
\end{enumerate}

\[
 m _ {i, j} = \operatorname{Card} \left\{f \in \mathrm{F} \mid o (f) = i \text {and} t (f) = j \right\}.
\]

The matrix \(M_{\mathrm{Q}}\) is called the adjacency matrix of Q.

\begin{enumerate}
\def\labelenumi{\alph{enumi})}
\item
  Given quivers Q and Q', give a necessary and sufficient condition on their adjacency matrices for Q and Q' to be isomorphic.
\item
  Express in terms of the adjacency matrix \(M_{Q}\) the number of paths of length n in Q whose origin and term are given.
\item
  Prove that the adjacency matrix of the underlying quiver of a graph is symmetric.
\end{enumerate}

\begin{enumerate}
\def\labelenumi{\arabic{enumi})}
\setcounter{enumi}{3}
\tightlist
\item
  Let \(Q = (S, F, o, t)\) be a quiver. We denote by I the sum set of the family \((S, F)\); we identify S and F with subsets of I. Let k be a commutative ring. The k-algebra of the quiver Q is called the quotient \(A_{Q}\) of the free associative k-algebra on the set I by the two-sided ideal generated by the elements \(X_{s}^{2} - X_{s}\) (for \(s \in S\)), \(X_{u}X_{v} - X_{v}X_{u}\) (for \(u, v \in S\)), \(X_{f}X_{t(f)} - X_{f}\) and \(X_{o(f)}X_{f} - X_{f}\) (for \(f \in F\)). For \(i \in I\), denote by \(x_{i}\) the image of \(X_{i}\) in the algebra \(A_{Q}\).
\end{enumerate}

\begin{enumerate}
\def\labelenumi{\alph{enumi})}
\item
  Show that one has \(\sum_{s\in S}x_s = 1\)
\item
  Prove that we have \(x_{f}x_{g}=0\) if f and g are arrows of Q that are not composable.
\item
  For every path \(c = (a_{0}, f_{1}, \ldots, f_{n}, a_{n})\) in Q, we set \(x_{c} = x_{a_{0}} x_{f_{1}} \ldots x_{f_{n}}\). Let c and \(c'\) be paths in Q; show that we have \(x_{c} x_{c'} = x_{cc'}\) if c and \(c'\) are composable, and \(x_{c} = x_{c'}\) otherwise.
\item
  Let C be the set of paths in the quiver Q; prove that the family \((x_{c})_{c\in C}\) is a basis of the k-module \(A_{Q}\).
\item
  Let \(J_{Q}\) be the two-sided ideal of \(A_{Q}\) generated by the elements of the form \(x_{f}\), for \(f \in F\). Prove that the algebra \(A_{Q}/J_{Q}\) is isomorphic to \(k^{(S)}\).
\item
  For the k-algebra \(A_{Q}\) to be a finitely generated k-module, it is necessary and sufficient that the set of vertices of Q be finite and that every loop in Q have length zero. The ideal \(J_{Q}\) is then nilpotent.
\end{enumerate}

We retain these notations for the rest of the exercises of §1. When we speak of an \(A_{Q}\)-module (without further specification), it is a right module. If A is a ring, an A-module M (without further specification) will be a right A-module; if a is an element of A, we denote by \(a_{M}\) or \((a)_{M}\) the homothety \(x \mapsto xa\) of M.

\begin{enumerate}
\def\labelenumi{\arabic{enumi})}
\setcounter{enumi}{4}
\tightlist
\item
  \begin{enumerate}
  \def\labelenumii{\alph{enumii})}
  \tightlist
  \item
    Prove that the quadruple \(Q^{\circ} = (S, F, t, o)\) is a quiver. It is called the opposite quiver of Q.
  \end{enumerate}
\end{enumerate}

\begin{enumerate}
\def\labelenumi{\alph{enumi})}
\setcounter{enumi}{1}
\tightlist
\item
  The algebra \(A_{Q^{\circ}}\) is isomorphic to the opposite algebra of the algebra \(A_{Q}\) .
\end{enumerate}

\begin{enumerate}
\def\labelenumi{\arabic{enumi})}
\setcounter{enumi}{5}
\tightlist
\item
  \begin{enumerate}
  \def\labelenumii{\alph{enumii})}
  \tightlist
  \item
    Let \((M_s)_{s\in\operatorname{Som}(Q)}\) be a family of \(k\)-modules and let \((u_f)_{f\in\operatorname{Fl}(Q)}\) be a family where, for each \(f\in\operatorname{Fl}(Q)\), \(u_f\) is a linear map from \(M_{o(f)}\) to \(M_{t(f)}\). Let M be the direct sum of the family \((M_s)\); for \(s\in S\), denote by \(j_s\) the canonical injection of \(M_s\) into M and \(p_s\) the canonical projection of M onto \(M_s\). Prove that there exists a unique structure of \(A_Q\)-module on M such that one has \((x_s)_M = j_s \circ p_s\) and \((x_f)_M = j_{t(f)} \circ u_f \circ p_{o(f)}\) for all \(s \in \operatorname{Som}(Q)\) and every \(f \in \operatorname{Fl}(Q)\).
  \end{enumerate}
\end{enumerate}

Conversely, every \(A_{Q}\) -module is of this form.

\begin{enumerate}
\def\labelenumi{\alph{enumi})}
\setcounter{enumi}{1}
\tightlist
\item
  If \(\Lambda\) is a \(k\)-module and \(s \in \operatorname{Som}(\mathbf{Q})\), denote by \(\Lambda(s)\) the \(\mathrm{A}_{\mathrm{Q}}\)-module associated with the family \((\mathrm{M}_i)_{i \in \operatorname{Som}(\mathbf{Q})}\) of \(k\)-modules given by \(\mathrm{M}_i = \Lambda\) if \(i = s\) and \(\mathrm{M}_i = 0\) otherwise, and with the family \((u_f)_{f \in \operatorname{Fl}(\mathbf{Q})}\) of linear maps, where \(u_f = 0\) for all \(f \in \operatorname{Fl}(\mathbf{Q})\).
\end{enumerate}

Let \(\Lambda\) and \(\Lambda'\) be nonzero \(k\)-modules. Let \(s\) and \(s'\) be vertices of Q. Prove that \(\Lambda(s)\) is isomorphic to \(\Lambda'(s')\) if and only if \(\Lambda\) and \(\Lambda'\) are isomorphic as \(k\)-modules and \(s = s'\).

\begin{enumerate}
\def\labelenumi{\alph{enumi})}
\setcounter{enumi}{2}
\item
  If \(\Lambda\) is a simple \(k\)-module, then \(\Lambda(s)\) is a simple \(\mathrm{A}_{\mathrm{Q}}\)-module for every vertex \(s\).
\item
  Assume further that every loop in Q has length zero. Prove that every simple \(A_Q\)-module is of the form \(\Lambda(s)\), where \(\Lambda\) is a simple \(k\)-module and s is a vertex of Q.
\end{enumerate}

\begin{enumerate}
\def\labelenumi{\arabic{enumi})}
\setcounter{enumi}{6}
\tightlist
\item
  Assume that \(k\) is an algebraically closed field. For every \(s \in \operatorname{Som}(Q)\), we set \(\mathrm{P}(s) = x_s\mathrm{A}_{\mathrm{Q}}\).
\end{enumerate}

\begin{enumerate}
\def\labelenumi{\alph{enumi})}
\item
  Prove that \(\mathrm{P}(s)\) is a projective and indecomposable \(A_{Q}\)-module, and that its socle is isomorphic to the \(A_{Q}\)-module \(k(s)\). Prove also that \(\mathrm{P}(s)/\mathrm{P}(s)\mathrm{J}_{\mathrm{Q}}\) is isomorphic to \(k(s)\).
\item
  Assume that every loop in Q has length zero. Prove that, for every \(s \in \operatorname{Som}(Q)\), the canonical homomorphism from \(k\) to \(\operatorname{End}(\mathrm{P}(s))\) is an isomorphism. Prove that, for every projective and indecomposable \(A_Q\)-module P, there exists a unique vertex s of Q such that P is isomorphic to \(\mathrm{P}(s)\).
\end{enumerate}

\begin{enumerate}
\def\labelenumi{\arabic{enumi})}
\setcounter{enumi}{7}
\tightlist
\item
  For \(s \in \operatorname{Som}(Q)\) we set \(\mathrm{P}(s) = x_s\mathrm{A}_{\mathrm{Q}}\) . Let M be a \(\mathrm{A}_{\mathrm{Q}}\) -module.
\end{enumerate}

\begin{enumerate}
\def\labelenumi{\alph{enumi})}
\tightlist
\item
  Let \(f \in \mathrm{Fl}(\mathbf{Q})\), let \(u = o(f)\) and \(v = t(f)\). Prove that there exists a unique pair \((\varphi_f', \varphi_f'')\) of homomorphisms of \(\mathrm{A}_{\mathrm{Q}}\)-modules
\end{enumerate}

\[
\varphi_ {f} ^ {\prime} \colon \mathrm{M} x _ {u} \otimes_ {k} \mathrm{P} (v) \to \mathrm{M} x _ {u} \otimes_ {k} \mathrm{P} (u), \quad \varphi_ {f} ^ {\prime \prime} \colon \mathrm{M} x _ {u} \otimes_ {k} \mathrm{P} (v) \to \mathrm{M} x _ {v} \otimes_ {k} \mathrm{P} (v),
\]

 such that \(\varphi_f'(m \otimes p) = m \otimes x_f p\) and \(\varphi_f''(m \otimes p) = mx_f \otimes p\) for all \(m \in \mathrm{M}x_u\) and every \(p \in \mathrm{P}(v)\).

\begin{enumerate}
\def\labelenumi{\alph{enumi})}
\setcounter{enumi}{1}
\item
  Let \(s \in \text{Som}(Q)\). Prove that there exists a unique homomorphism of \(A_{Q}\)-modules \(\psi_{s}: Mx_{s} \otimes_{k} P(s) \to M\) such that \(\psi_{s}(m \otimes p) = mp\) for \(p \in P(s)\) and \(m \in Mx_{s}\).
\item
  Set \(\varphi = \bigoplus \varphi_f' - \bigoplus \varphi_f''\) and \(\psi = \sum \psi_s\). Show that one has an exact sequence of \(\mathrm{A}_{\mathrm{Q}}\)-modules
\end{enumerate}

\[
0 \to \bigoplus_ {f \in \operatorname{Fl} (\mathrm{Q})} \operatorname{M} x _ {o (f)} \otimes_ {k} \operatorname{P} (t (f)) \xrightarrow {\varphi} \bigoplus_ {s \in \operatorname{Som} (\mathrm{Q})} \operatorname{M} x _ {s} \otimes_ {k} \operatorname{P} (s) \xrightarrow {\psi} \operatorname{M} \to 0.
\]

\(d\)) Show that every right ideal of the algebra \(\mathrm{A_Q}\) is projective.

\begin{enumerate}
\def\labelenumi{\arabic{enumi})}
\setcounter{enumi}{8}
\tightlist
\item
  \begin{enumerate}
  \def\labelenumii{\alph{enumii})}
  \tightlist
  \item
    Let M and N be \(A_Q\)-modules; for \(s \in \operatorname{Som}(Q)\), set \(M_s=Mx_s\) and \(N_s=Nx_s\). For \(\varphi = (\varphi_s)_{s \in \operatorname{Som}(Q)}\) and \(f \in \operatorname{Fl}(Q)\), define a map \(\theta_f(\varphi): Mx_{o(f)} \to Nx_{t(f)}\) by setting
  \end{enumerate}
\end{enumerate}

\[
\theta_ {f} (\varphi) (m) = \varphi_ {o (f)} (m) x _ {f} - \varphi_ {t (f)} (m x _ {f}).
\]

Prove that the map

\[
\theta \colon \prod_ {s \in \operatorname{Som} (\mathrm{Q})} \operatorname{Hom} _ {k} \left(\mathrm{M} _ {s}, \mathrm{N} _ {s}\right)\rightarrow \prod_ {f \in \operatorname{Fl} (\mathrm{Q})} \operatorname{Hom} \left(\mathrm{M} _ {o (f)}, \mathrm{N} _ {t (f)}\right)
\]

 given by \(\varphi\mapsto(\theta_{f}(\varphi))_{f\in\mathrm{Fl}(Q)}\) is linear.

\begin{enumerate}
\def\labelenumi{\alph{enumi})}
\setcounter{enumi}{1}
\item
  Prove that \(\operatorname{Ker}(\theta)\) is isomorphic to \(\operatorname{Hom}_{\mathrm{A}_{\mathrm{Q}}}(\mathrm{M},\mathrm{N})\) and that \(\operatorname{Coker}(\theta)\) is isomorphic to \(\operatorname{Ext}_{\mathrm{A}_{\mathrm{Q}}}^{1}(\mathrm{M},\mathrm{N})\). Prove that \(\operatorname{Ext}_{\mathrm{A}_{\mathrm{Q}}}^{p}(\mathrm{M},\mathrm{N}) = 0\) for every integer \(p \geqslant 2\).
\item
  Suppose that k is a commutative field and that the quiver Q is finite. If M and N are finite-dimensional k-vector spaces, then \(\operatorname{Ext}_{\mathrm{A}_{\mathrm{Q}}}^{p}(\mathrm{M},\mathrm{N})\) is a finite-dimensional k-vector space for every integer \(p \geqslant 0\), is zero for \(p \geqslant 2\), and
\end{enumerate}

\[
\begin{array}{l} \dim_ {k} (\mathrm{Hom} _ {\mathrm{A} _ {\mathrm{Q}}} (\mathrm{M}, \mathrm{N})) - \dim_ {k} (\mathrm{Ext} _ {\mathrm{A} _ {\mathrm{Q}}} ^ {1} (\mathrm{M}, \mathrm{N})) \\ = \sum_ {s \in \mathrm{Som} (\mathrm{Q})} \dim (\mathrm{M} _ {s}) \dim (\mathrm{N} _ {s}) - \sum_ {f \in \mathrm{Fl} (\mathrm{Q})} \dim (\mathrm{M} _ {o (f)}) \dim (\mathrm{N} _ {t (f)}). \end{array}
\]

\(d\)) Assume that \(k\) is a commutative field. Prove that for every pair \((u,v)\) of vertices of Q, the dimension of the vector space \(\mathrm{Ext}_{\mathrm{A_Q}}^1 (k(u),k(v))\) is equal to the cardinality of the set of arrows of Q with origin \(v\) and target \(u\).

\begin{enumerate}
\def\labelenumi{\arabic{enumi})}
\setcounter{enumi}{9}
\item
  We assume that \(k\) is an algebraically closed field. Let Q and Q' be finite quivers in which every loop has length zero. Prove that the algebra A \(_{Q'}\) is Morita equivalent (A, VIII, §6) to the algebra A \(_{Q}\) if and only if the quivers Q and Q' are isomorphic. (Take as modules M, N the \(k\) -simple modules \(k(s)\) , for \(s \in \text{Som}(Q)\) .)
\end{enumerate}

\exercisegroup

\begin{enumerate}
\def\labelenumi{\arabic{enumi})}
\item
  Let G be a graph and let \(Q = (S, F, o, t)\) its underlying quiver.
\end{enumerate}

\begin{enumerate}
\def\labelenumi{\alph{enumi})}
\item
  Show that one has \(d_{+}(s) = d_{-}(s)\) for all \(s \in S\). This integer is then called the degree of G at s and denoted \(d(s)\).
\item
  Assume that G is finite. For there to exist a closed path \((a_{0}, f_{1}, \ldots, f_{n}, a_{n})\) in G such that, for every edge \(\varphi\) of G, there exists a unique integer n such that \(\varphi = \{f_{n}, \overline{f_{n}}\}\), it is necessary and sufficient that G be connected and that \(d(s)\) be even for all \(s \in \text{Som}(G)\).
\end{enumerate}

\begin{enumerate}
\def\labelenumi{\arabic{enumi})}
\setcounter{enumi}{1}
\tightlist
\item
  Let \(\mathrm{G} = (\mathrm{S}, \mathrm{F}, o, t)\) be a quiver. We say that G is a bipartite quiver if there exists a partition \(P = \{S_{1}, S_{2}\}\) of S into subsets such that, for every arrow \(f \in F\), the origin \(o(f)\) and the term \(t(f)\) are in different parts of the partition P.
\end{enumerate}

\begin{enumerate}
\def\labelenumi{\alph{enumi})}
\item
  Assume that G is connected. Prove that there exists at most one partition P as above.
\item
  Prove that a graph G is bipartite if and only if every cycle in G has even length.
\end{enumerate}

\begin{enumerate}
\def\labelenumi{\arabic{enumi})}
\setcounter{enumi}{2}
\item
  Let G be a connected, locally finite graph whose vertex set is infinite. Prove that there exists a composable sequence \((f_{n})_{n\in\mathbf{N}}\) of arrows of G whose origins are pairwise distinct (König's theorem).
\item
  Let G be a finite graph and let \(Q = (S, F, o, t)\) be its underlying quiver. We endow the space \(C^{S}\) with the unique Hermitian space structure for which its canonical basis is orthonormal. Let L: \(C^{S} \to C^{S}\) be the map given by \(\mathrm{L}(u)(x) = \sum_{o(f)=x} (u(t(f)) - u(x))\) for \(x \in S\) and \(u \in C^{S}\).
\end{enumerate}

\begin{enumerate}
\def\labelenumi{\alph{enumi})}
\item
  Prove that the endomorphism L of \(C^{S}\) is positive (EVT, V, p. 45, def. 6).
\item
  Prove that the kernel of L consists of the functions \(u \in C^{S}\) which are constant on each connected component of G.
\item
  Let \(F_{1}\) be an orientation of G. Let \(E = (e_{s,f})\) be the matrix of type \((\mathrm{S}, \mathrm{F}_{1})\) given by
\end{enumerate}

\[
 e _ {s, f} = \left\{ \begin{array}{l l} 1 & \text {if } o (f) = s \text {and } t (f) \neq s; \\ - 1 & \text {if } o (f) \neq s \text {and } t (f) = s; \\ 0 & \text {otherwise. } \end{array} \right.
\]

Prove that the matrix \(\Lambda\) of L in the canonical basis of \(\mathbf{C}^{\mathrm{S}}\) is given by \(\Lambda = E\cdot {}^{t}E\) .

\(d)\) Assume that G is connected and that its set of vertices is nonempty; let \(x\) be a vertex of G and let \(\mathrm{S}' = \mathrm{S} - \{x\}\). Let T be a subset of F of cardinality \(\mathrm{Card}(\mathrm{S}')\). Prove that T is the set of arrows of a maximal oriented subtree of G (equipped with the orientation \(\mathrm{F}_1\)) if and only if the rank of the matrix \(E_{\mathrm{T}}'\) of type \((\mathrm{S}',\mathrm{T})\) deduced from \(E\) is equal to \(\mathrm{Card}(\mathrm{S}')\). From this, deduce that the determinant of the matrix \(A_{\mathrm{S}',\mathrm{S}'}\) is equal to the number of maximal oriented subtrees of G. (Use exercise 6 of A, III, §8, p. 192.)

\begin{enumerate}
\def\labelenumi{\alph{enumi})}
\setcounter{enumi}{4}
\tightlist
\item
  Let W be the orthogonal complement of \(\mathrm{Ker}(\mathrm{L})\). Prove that W is a subspace of \(\mathbf{C}^{\mathrm{S}}\) that is stable under W and that \(\det (\mathrm{L}|_{\mathrm{W}})\) is equal to the number of maximal forests of the graph G (Kirchhoff's theorem).
\end{enumerate}

\begin{enumerate}
\def\labelenumi{\arabic{enumi})}
\setcounter{enumi}{4}
\tightlist
\item
  \begin{enumerate}
  \def\labelenumii{\alph{enumii})}
  \tightlist
  \item
    Let n be a natural integer \(\geqslant 1\) and let G be a graph whose set of vertices is the set \(\mathbf{Z}/n\mathbf{Z}\) and such that two vertices i and j are joined by an arrow if and only if \(i-j=\pm1\pmod n\). Describe explicitly the set of maximal forests of G.
  \end{enumerate}
\end{enumerate}

\begin{enumerate}
\def\labelenumi{\alph{enumi})}
\setcounter{enumi}{1}
\item
  Let G be a graph and let S be the set of its vertices. Suppose that, for every pair \((x,y)\) of vertices of G, the set of arrows of G with origin x and terminus y has cardinality 1 if \(x \neq y\) and 0 otherwise ("complete graph"). Compute the cardinality of the set of maximal forests of G.
\item
  Let G be a graph and let S be the set of its vertices. Let \((\mathrm{S}_{1},\mathrm{S}_{2})\) be a partition of S; suppose that, for every pair \((x,y)\) of vertices of G, the set of arrows of G with origin x and terminus y has cardinality 1 if \((x,y)\in\mathrm{S}_{1}\times\mathrm{S}_{2}\) and 0 otherwise ("complete bipartite graph"). Compute the cardinality of the set of maximal forests of G.
\end{enumerate}

\begin{enumerate}
\def\labelenumi{\arabic{enumi})}
\setcounter{enumi}{5}
\tightlist
\item
  Let \(c\) be a natural integer.
\end{enumerate}

\begin{enumerate}
\def\labelenumi{\alph{enumi})}
\tightlist
\item
  There exists a unique family of maps \(\mathrm{R}_{c}\colon(\mathbf{N}^{*})^{c}\to\mathbf{N}^{*}\), for \(c\geqslant2\), satisfying the following properties:
\end{enumerate}

\begin{enumerate}
\def\labelenumi{(\roman{enumi})}
\item
  If there exists i such that \(n_{i}=1\), then \(\mathrm{R}_{c}(n_{1},\ldots,n_{c})=1\);
\item
  We have \(\mathrm{R}_{2}(m,n)=\mathrm{R}_{2}(m-1,n)+\mathrm{R}_{2}(m,n-1)\) if m and n are integers \(\geqslant2\);
\item
  If \(c \geqslant 3\), we have \(\mathrm{R}_c(n_1, \ldots, n_c) = \mathrm{R}_{c-1}(n_1, \ldots, n_{c-2}, \mathrm{R}_2(n_{c-1}, n_c))\) for all \((n_1, \ldots, n_c) \in \mathbf{N}^c\).
\end{enumerate}

\begin{enumerate}
\def\labelenumi{\alph{enumi})}
\setcounter{enumi}{1}
\item
  We have \(\mathrm{R}_2(m,n) = \binom{m + n - 2}{m - 1}\) for every pair \((m,n)\) of integers \(\geqslant 1\).
\item
  For every finite sequence \((n_{1},\ldots,n_{c})\) of natural integers \(\geqslant1\) and every natural integer \(R\geqslant1\), we say that the property \(\boldsymbol{R}(n_{1},\ldots,n_{c};\mathrm{R})\) is satisfied if, for every complete graph whose set of vertices S has cardinality \(\geqslant R\) and every partition \(P=(F_{1},\ldots,F_{c})\) of the set F of its arrows, there exists an integer \(i\in\{1,\ldots,c\}\) and a subset A of S of cardinality \(n_{i}\) such that every arrow of G connecting two elements of A belongs to \(F_{i}\).
\end{enumerate}

Prove the property \(\boldsymbol{R}(n_{1},\ldots,n_{c};\mathrm{R}_{c}(n_{1},\ldots,n_{c}))\) .

\begin{enumerate}
\def\labelenumi{\alph{enumi})}
\setcounter{enumi}{3}
\item
  We denote by \(\mathrm{R}(n_{1},\ldots,n_{c})\) the smallest integer \(R\geqslant1\) such that the property \(\boldsymbol{R}(n_{1},\ldots,n_{c};\mathrm{R})\) is satisfied (“Ramsey numbers”).
\item
  We have \(R(n,2)=n\) for every integer \(n\geqslant2\); we have \(R(3,3)=6\).
\item
  Prove that R(4,3) = 9.
\end{enumerate}

Compute R(5, 5); compute R(6, 6).

\begin{enumerate}
\def\labelenumi{\arabic{enumi})}
\setcounter{enumi}{6}
\tightlist
\item
  Let G be a group acting on the right on a set X, and let A be a subset of G. The Schreier graph \(\mathcal{G}_{\mathrm{G}}(\mathrm{X},\mathrm{A})\) is the quiver \((\mathrm{S},\mathrm{F},o,t)\) defined by:
\end{enumerate}

- The set of its vertices is S = X;

- The set of its arrows is F = X × A;

- The maps \(o\) and \(t\) are given by \(o(x, a) = x\) and \(t(x, a) = x \cdot a\) for all \((x, a) \in \mathrm{X} \times \mathrm{A}\) .

  When X = G, on which G acts by right multiplication, we denote this quiver by \(\mathcal{G}(G,A)\) and call it the Cayley quiver of G with respect to A.

We call the graph associated with this quiver a Schreier graph (resp. Cayley graph).

\begin{enumerate}
\def\labelenumi{\alph{enumi})}
\item
  Prove that the quiver \(\mathcal{G}_{\mathrm{G}}(\mathrm{X},\mathrm{A})\) is connected if and only if the subgroup H of G generated by A has at most one orbit on X.
\item
  Assume that A generates G. Show that the Cayley graph \(\mathcal{G}(\mathrm{G},\mathrm{A})\) is connected and that it is bipartite if and only if there exists a group homomorphism \(\varphi \colon \mathrm{G}\to \mathbf{Z} / 2\mathbf{Z}\) such that \(\varphi (a) = 1\) for all \(a\in \mathrm{A}\) .
\item
  Show that there exists a unique action of G on the Cayley graph \(\mathcal{G}(\mathrm{G},\mathrm{A})\) such that the action of G induced on the vertices of this graph is the action of G on itself by left multiplication.
\end{enumerate}

\begin{enumerate}
\def\labelenumi{\arabic{enumi})}
\setcounter{enumi}{7}
\tightlist
\item
  A cycle \((a_0, f_1, \ldots, f_n, a_n)\) in a graph G is said to be Hamiltonian if, for every vertex \(s\) of G, there exists a unique integer \(m \in \{1, \ldots, n\}\) such that \(a_m = s\). A graph is said to be Hamiltonian if it has a Hamiltonian cycle in G.
\end{enumerate}

\begin{enumerate}
\def\labelenumi{\alph{enumi})}
\item
  Let S be a finite set of cardinality \(\geqslant 3\) and let \(K_{S}\) be a complete graph with vertex set S. Prove that the graph \(K_{S}\) is Hamiltonian. More precisely, prove that for every orientation A of G, there exists a Hamiltonian cycle in \(K_{S}\) whose every arrow belongs to A.
\item
  We denote by \(G_{S}\) the set of subgraphs of \(K_{S}\) whose set of vertices is equal to S. Let \(\preccurlyeq\) be the coarsest order relation on the set \(G_{S}\) for which \(G \preccurlyeq G'\) if there exist elements u, v of S such that \(\deg(u) + \deg(v) \geqslant \text{Card}(S)\) and such that \(\text{Fl}(G')\) is the union of \(\text{Fl}(G)\) and the two arrows of \(K_{S}\) with endpoints u and v.
\end{enumerate}

  Prove that, for every \(G \in G_{S}\), the set of elements H of \(G_{S}\) such that \(G \preccurlyeq H\) has a unique maximal element; it is denoted \(c(G)\).

\begin{enumerate}
\def\labelenumi{\alph{enumi})}
\setcounter{enumi}{2}
\item
  Let G be a subgraph of \(K_{S}\) with vertex set S. Prove that the graph G is Hamiltonian if and only if the same is true of the graph \(c(\mathrm{G})\) .
\item
  Let G be a subgraph of \(K_{S}\) whose vertex set is S and whose degree at each vertex is \(\geqslant\) Card(S)/2. Prove that \(c(\mathrm{G})\) is equal to \(K_{S}\) . Deduce that G is a Hamiltonian graph (theorem of G. Dirac).
\end{enumerate}

\begin{enumerate}
\def\labelenumi{\arabic{enumi})}
\setcounter{enumi}{8}
\tightlist
\item
  Let (W, S) be a Coxeter system (LIE, IV, §1, p. 11, def. 3); suppose that the group W is finite.
\end{enumerate}

\begin{enumerate}
\def\labelenumi{\alph{enumi})}
\item
  Suppose that (W, S) is irreducible and has at most two vertices. Prove that the Cayley graph \(\mathcal{G}(\mathrm{W},\mathrm{S})\) is Hamiltonian.
\item
  Suppose that (W, S) is irreducible and has at least three vertices. Let \(s \in S\) be a terminal vertex (LIE IV, appendix) of the Coxeter graph associated with (W, S). Set \(S' = S - \{s\}\) and denote by \(W'\) the subgroup of W generated by \(S'\). Let \(\Gamma\) be the graph whose set of vertices is the set of left cosets modulo \(W'\) in W and such that two distinct vertices A and B are joined by an edge exactly if and only if \(\cap\) is not empty. Using a maximal subtree of the graph \(\Gamma\), prove that if the Cayley graph \(\mathcal{G}(W', S')\) is Hamiltonian, then so is the graph \(\mathcal{G}(W, S)\).
\item
  Prove that the Cayley graph \(\mathcal{G}(\mathrm{W},\mathrm{S})\) is Hamiltonian (Conway-Sloane-Wilks theorem).
\end{enumerate}

\begin{enumerate}
\def\labelenumi{\arabic{enumi})}
\setcounter{enumi}{9}
\tightlist
\item
  Let G be a connected graph whose set of vertices is S. For \(x, y \in S\), we set
\end{enumerate}

\(d_{\mathrm{G}}(x,y) = \inf \{\ell \in \mathbf{N}\colon \text{il existe un chemin de longueur} \ell \text{ reliant } x \text{ à } y\}.\)

\begin{enumerate}
\def\labelenumi{\alph{enumi})}
\item
  Prove that \(d_{\mathrm{G}}\) is a distance on the set S.
\item
  Let \(G_{1}\) and \(G_{2}\) be graphs and \(\varphi\colon G_{1}\to G_{2}\) a graph morphism. Show that \(d_{\mathrm{G}_{2}}(\varphi(x),\varphi(y))\leqslant d_{\mathrm{G}_{1}}(x,y)\) for every pair \((x,y)\) of vertices of \(G_{1}\) .
\end{enumerate}

\begin{enumerate}
\def\labelenumi{\arabic{enumi})}
\setcounter{enumi}{10}
\tightlist
\item
  Let G be a connected graph whose vertex set S is nonempty. We call the diameter of G the diameter \(\operatorname{diam}(G)\) of the metric space (S, \(d_{\mathrm{G}}\)).
\end{enumerate}

\begin{enumerate}
\def\labelenumi{\alph{enumi})}
\item
  Prove that \(\text{diam}(G) \leqslant \text{Card}(S)-1\). For which graphs is this inequality an equality?
\item
  Let v be an integer \(\geqslant 2\) such that every vertex of G has degree at most v. Prove that \(\text{Card}(S) \leqslant v^{\text{diam}(G)}\). For which graphs is this inequality an equality?
\item
  Let n be an integer such that \(n \geqslant 2\). Let \(T_{n}\) be the subset of \(S_{n}\) formed by the transpositions with support \(\{m, m + 1\}\), for \(m \in \{1, \ldots, n - 1\}\). Prove that the diameter of the Cayley graph \(\mathcal{G}(\mathfrak{S}_{n}, \mathrm{T}_{n})\) is equal to \(n(n - 1)/2\).
\item
  Let n be an integer such that \(n \geqslant 2\). Let \(\tau\) be the transposition with support \(\{1,2\}\) and \(\sigma\) be the cycle \((1,2,\ldots,n) \mapsto (2,3,\ldots,n,1)\) in the group \(S_n\). Let \(G_n\) be the Cayley graph \(\mathcal{G}(\mathfrak{S}_n,\{\sigma,\tau\})\). Prove that \(n(n-1)/6 \leqslant \text{diam}(G_n) \leqslant 3n(n-1)/2\). (Let T be the set of sequences \((x,y,z)\) of elements of \(\{1,\ldots,n\}\) such that \(x \textless{} y \textless{} z\). If \(t=(x,y,z)\in\text{T}\), we say that a permutation \(\lambda\in S_n\) preserves the cyclic order of t if one has \(\lambda(x)<\lambda(y)<\lambda(z)\), or \(\lambda(y)<\lambda(z)<\lambda(x)\), or \(\lambda(z)<\lambda(x)<\lambda(y)\). Observe that the permutation \(\sigma\) preserves the cyclic order of every element of T, and that the permutation \(\tau\) preserves the cyclic order of every element of T that is not of the form \((1,2,x)\). Deduce that the distance from the identity permutation to the permutation \((1,\ldots,n) \mapsto (n,n-1,\ldots,1)\) is at least equal to \(n(n-1)/6\).
\end{enumerate}

\begin{enumerate}
\def\labelenumi{\arabic{enumi})}
\setcounter{enumi}{11}
\tightlist
\item
  Let \(n\) be an integer \(\geqslant 2\) and let \(p\) be a prime number \(>n/2\). Let H be a subgroup of \(\mathfrak{S}_n\) that acts transitively on \(\{1, \ldots, n\}\), contains a transposition and a cycle of order \(p\). Let G be a graph such that
\end{enumerate}

- The set of its vertices is \(S = \{1, \ldots, n\}\) ;

- The set F of its arrows is the set of pairs \((i,j)\in\mathrm{S}\times\mathrm{S}\) such that the transposition \(\tau_{i,j}\) belongs to H;

- We have \(o((i,j)) = i\) , \(t((i,j)) = j\) and \(\overline{(i,j)} = (j,i)\) for all \((i,j) \in \mathrm{F}\) .

\begin{enumerate}
\def\labelenumi{\alph{enumi})}
\item
  Prove that every connected component of G is a complete graph.
\item
  Prove that if G is connected, then \(H = S_{n}\) .
\item
  Prove that there exists a group homomorphism from \(\mathfrak{S}_n\) to \(\operatorname{Aut}(\mathbf{G})\) such that, for every \(\sigma \in \mathfrak{S}_n\), the map \(\operatorname{Som}(\sigma)\) is equal to the permutation \(\sigma\). Then show that \(\mathfrak{S}_n\) acts transitively on the set of connected components of \(\mathbf{G}\), and that these are pairwise isomorphic.
\end{enumerate}

\(d)\) Let \(\sigma\in\mathrm{H}\) be a cycle of order \(p\). Prove that \(\sigma\) stabilizes each connected component of G. Conclude that \(\mathrm{H}=\mathfrak{S}_{n}\).

\exercisegroup

\begin{enumerate}
\def\labelenumi{\arabic{enumi})}
\tightlist
\item
  Let \((\mathrm{C}_{i})_{i\in\mathrm{I}}\) be a family of categories. Let C be the sum quiver of the family of underlying quivers of the categories \(C_{i}\).
\end{enumerate}

\begin{enumerate}
\def\labelenumi{\alph{enumi})}
\item
  Show that the quiver C, equipped with the composition law induced by the composition laws of the categories \(C_{i}\), is a category. We call it the sum category of the family \(C_{i}\).
\item
  Prove that the canonical morphism from \(C_{i}\) to C is a functor.
\item
  If the \(C_{i}\) are groupoids, prove that the category C is a groupoid.
\end{enumerate}

\begin{enumerate}
\def\labelenumi{\arabic{enumi})}
\setcounter{enumi}{1}
\tightlist
\item
  Let \((\mathrm{C}_{i})_{i\in\mathrm{I}}\) be a family of categories. Let C be the product of the underlying quivers of the categories \(C_{i}\).
\end{enumerate}

\begin{enumerate}
\def\labelenumi{\alph{enumi})}
\item
  Show that the quiver C, equipped with the composition law induced by the composition laws of the categories \(C_{i}\), is a category. It is called the product category of the family \((\mathrm{C}_{i})\).
\item
  Prove that the canonical morphism \(pr_{i}: C \rightarrow C_{i}\) is a functor.
\item
  If the \(C_{i}\) are groupoids, prove that the category C is a groupoid.
\end{enumerate}

\begin{enumerate}
\def\labelenumi{\arabic{enumi})}
\setcounter{enumi}{2}
\tightlist
\item
  Let C be a category; denote its composition law by m. Let C° be the opposite directed graph equipped with the composition law given by \(m^{\circ}(f,g)=m(g,f)\) for every pair \((f,g)\) of composable arrows of C.
\end{enumerate}

\begin{enumerate}
\def\labelenumi{\alph{enumi})}
\item
  Prove that C° is a category. It is called the opposite category of C.
\item
  If C is a groupoid, prove that C° is a groupoid isomorphic to C.
\end{enumerate}

\begin{enumerate}
\def\labelenumi{\arabic{enumi})}
\setcounter{enumi}{3}
\tightlist
\item
  We define the mass \(\mu(\mathrm{G})\) of a groupoid G as the supremum of the sums \(\sum_{a\in\mathrm{A}}\mathrm{Card}(\mathrm{G}_{a})^{-1}\), where A ranges over the set of subsets of \(\mathrm{Som}(\mathrm{G})\) which meet each orbit of G in at most one point. We say that G has finite mass if \(\mu(\mathrm{G})\in\mathbf{R}\); we say that G is essentially finite if the set of its orbits is finite and if the isotropy group at each of its points is finite.
\end{enumerate}

\begin{enumerate}
\def\labelenumi{\alph{enumi})}
\item
  Show that an essentially finite groupoid has finite mass.
\item
  Let G be the groupoid associated with a finite group \(\Gamma\). It is essentially finite, and we have \(\mu(\mathrm{G}) = 1/\mathrm{Card}(\Gamma)\).
\item
  Let \(\Gamma\) be a finite group, let X be a finite set equipped with a right action of \(\Gamma\), and let G be the associated groupoid. It is essentially finite, and we have \(\mu(\mathrm{G}) = \mathrm{Card}(\mathrm{X}) / \mathrm{Card}(\Gamma)\).
\end{enumerate}

\(d)\) Let \((\mathrm{G}_i)_{i\in \mathrm{I}}\) be a family of groupoids of finite mass. The sum groupoid and the product groupoid of the family \((\mathrm{G}_i)\) have mass \(\sum \mu (\mathrm{G}_i)\) and \(\prod \mu (\mathrm{G}_i)\), respectively.

\begin{enumerate}
\def\labelenumi{\alph{enumi})}
\setcounter{enumi}{4}
\tightlist
\item
  Let X be a set. We define a groupoid G whose set of vertices is the set of finite subsets of X and such that, for every pair \((Y_{1}, Y_{2})\) of finite subsets of X, the set of arrows with origin \(Y_{1}\) and term \(Y_{2}\) is the set of bijections from \(Y_{1}\) to \(Y_{2}\), with composition of arrows given by composition of maps. Show that G is essentially finite and compute its mass.
\end{enumerate}

\begin{enumerate}
\def\labelenumi{\arabic{enumi})}
\setcounter{enumi}{4}
\tightlist
\item
  \begin{enumerate}
  \def\labelenumii{\alph{enumii})}
  \tightlist
  \item
    Let \(\varphi\colon G\to G'\) be a groupoid morphism such that the map \(\operatorname{Som}(\varphi)\) is injective. Prove that \(\varphi(G)\) is a subgroupoid of \(G'\).
  \end{enumerate}
\end{enumerate}

\begin{enumerate}
\def\labelenumi{\alph{enumi})}
\setcounter{enumi}{1}
\tightlist
\item
  Let G be a nontransitive groupoid whose vertex set \(\operatorname{Som}(G)\) is the two-element set \(\{a,b\}\). Let G' be the groupoid associated with the free product \(G_a*G_b\). There exists a unique groupoid morphism \(\varphi:G\to G'\) such that the maps \(\varphi_a\) and \(\varphi_b\) are the canonical maps from \(G_a\) and \(G_b\) to \(G_a*G_b\). If \(G_a\) or \(G_b\) is not the trivial group, the image \(\varphi(G)\) of this groupoid morphism is not a subgroupoid of \(G'\).
\end{enumerate}

\begin{enumerate}
\def\labelenumi{\arabic{enumi})}
\setcounter{enumi}{5}
\tightlist
\item
  Let G and G' be graphs, and let \(\varphi\colon\mathrm{G}'\to\mathrm{G}\) be a graph morphism. We say that \((\mathrm{G}',\varphi)\) is a covering of G if, for every vertex \(a\in\mathrm{Som}(\mathrm{G}')\) and every arrow \(f\in\mathrm{Fl}(\mathrm{G})\) with origin \(\mathrm{Som}(\varphi)(a)\), there exists a unique arrow \(f'\in\mathrm{Fl}(\mathrm{G}')\) with origin \(a\) such that \(\mathrm{Fl}(\varphi)(f')=f\).
\end{enumerate}

We assume that \((\mathrm{G}^{\prime},\varphi)\) is a covering of G.

\begin{enumerate}
\def\labelenumi{\alph{enumi})}
\item
  Prove that there exists a unique action of \(\operatorname{Grp}(G)\) on \(\operatorname{Som}(G')\), relative to the map \(\operatorname{Som}(\varphi)\), such that \(x \cdot \operatorname{Fl}(\varphi)(f) = t(f)\) for every vertex \(x \in \operatorname{Som}(G')\) and every arrow \(f \in \operatorname{Fl}(G')\) with origin x.
\item
  Deduce that if x and y are vertices of G belonging to the same connected component of G, then \(\text{Card}(\text{Som}(\varphi)^{-1}(x)) = \text{Card}(\text{Som}(\varphi)^{-1}(y))\).
\end{enumerate}

\begin{enumerate}
\def\labelenumi{\arabic{enumi})}
\setcounter{enumi}{6}
\tightlist
\item
  Let G be a graph and let \((G_{1},\varphi_{1})\), \((G_{2},\varphi_{2})\) be coverings of the graph G.
\end{enumerate}

\begin{enumerate}
\def\labelenumi{\alph{enumi})}
\item
  Let \(\psi\colon\mathrm{G}_{1}\to\mathrm{G}_{2}\) be a graph morphism such that \(\varphi_{2}\circ\psi=\varphi_{1}\). Show that \(\psi(x\cdot\gamma)=\psi(x)\cdot\gamma\) for all \(x\in\mathrm{Som}(\mathrm{G}_{1})\) and every arrow \(\gamma\) of \(\varpi_{\mathrm{G}}\) with origin \(\varphi_{1}(x)\).
\item
  Let \(u \colon \operatorname{Som}(\mathrm{G}_1) \to \operatorname{Som}(\mathrm{G}_2)\) be a map such that \(u(x \cdot \gamma) = u(x) \cdot \gamma\) for all \(x \in \operatorname{Som}(\mathrm{G}_1)\) and every arrow \(\gamma\) of \(\varpi_{\mathrm{G}}\) with origin \(\varphi_1(x)\). Prove that there exists a unique morphism of graphs \(\psi \colon \mathrm{G}_1 \to \mathrm{G}_2\) such that \(\varphi_2 \circ \psi = \varphi_1\) and \(\operatorname{Som}(\psi) = u\).
\item
  Let X be a set and let \(p: X \to \text{Som}(G)\) be a map. Suppose that an operation of \(\varpi_{G}\) on X is given with respect to p. Prove that there exists a graph \(G'\) with vertex set X and a morphism of graphs \(\varphi: G' \to G\) such that \(\text{Som}(\varphi) = p\), such that \((G', \varphi)\) is a covering of G and the operation of \(\varpi_{G}\) on \(\text{Som}(G')\) coincides with the given operation.
\end{enumerate}

\begin{enumerate}
\def\labelenumi{\arabic{enumi})}
\setcounter{enumi}{7}
\tightlist
\item
  Let G be a connected graph and let \(a\) be a vertex of G.
\end{enumerate}

\begin{enumerate}
\def\labelenumi{\alph{enumi})}
\tightlist
\item
  Let \((\mathrm{G}_{1},\varphi_{1})\) and \((\mathrm{G}_{2},\varphi_{2})\) be coverings of the graph G. Let \(u\colon\operatorname{Som}(\mathrm{G}_{1})\to\operatorname{Som}(\mathrm{G}_{2})\) be a map such that \(u(x)\cdot\gamma=u(x\cdot\gamma)\)
\end{enumerate}

for all \(\gamma \in \varpi_{\mathrm{G},a}\) and every vertex \(x\) of \(\mathrm{G}_1\) such that \(\operatorname{Som}(\varphi_1)(x) = a\). Prove that there exists a unique morphism of graphs \(\psi: \mathrm{G}_1 \to \mathrm{G}_2\) such that \(\psi = \operatorname{Som}(u)\).

\begin{enumerate}
\def\labelenumi{\alph{enumi})}
\setcounter{enumi}{1}
\tightlist
\item
  Let X be a set equipped with an operation of the group \(\varpi_{\mathrm{G},a}\). Prove that there exists a covering \((\mathrm{G}',\varphi)\) of the graph G such that \(\operatorname{Som}(\varphi)^{-1}(a) = \mathrm{X}\) and such that the operation of \(\varpi_{\mathrm{G}}\) on \(\operatorname{Som}(\mathrm{G}')\) induces the given operation of \(\varpi_{\mathrm{G},a}\) on X.
\end{enumerate}

\begin{enumerate}
\def\labelenumi{\arabic{enumi})}
\setcounter{enumi}{8}
\tightlist
\item
  Let G and G' be graphs, and let \(\varphi\colon G'\to G\) be a graph morphism such that \((G',\varphi)\) is a covering of G.
\end{enumerate}

\begin{enumerate}
\def\labelenumi{\alph{enumi})}
\item
  Prove that the graph \(\varphi(\mathrm{Grp}(\mathrm{G}^{\prime}))\) is a subgroupoid of the groupoid \(\mathrm{Grp}(\mathrm{G})\) .
\item
  Prove that, for every vertex \(a\) of \(\mathbf{G}'\), the group homomorphism \(\pi_1(\varphi, a)\) is injective.
\item
  Suppose that G is connected and fix a vertex a of G. Let K be a subgroup of \(\varpi_{G,a}\). Prove that there exists a covering \((H,\psi)\) of the graph G and a vertex b of H such that \(\psi(b)=a\) and K is the image of the morphism \(\pi_{1}(\psi,a)\).
\item
  Prove that a subgroup of a free group is free (see also IV, p. 417, example 2).
\item
  Let n and q be integers; if F is a free group on n generators and K a subgroup of F of index q, prove that K is a free group on \(1 + q(n - 1)\) generators.
\end{enumerate}

\begin{enumerate}
\def\labelenumi{\arabic{enumi})}
\setcounter{enumi}{9}
\tightlist
\item
  Let F be a set, and let C be a subset of \(F \times F\) and \(m: C \to F\) a map. We make the following assumptions:
\end{enumerate}

\begin{enumerate}
\def\labelenumi{(\roman{enumi})}
\item
  The maps from C into F × F given by \((f,g)\mapsto(f,m(f,g))\) and \((f,g)\mapsto(m(f,g),g)\) are injective;
\item
  If \(f, g, h \in \mathrm{F}\) are such that the pairs \((f, g)\) and \((g, h)\) belong to C, then the pairs \((f, m(g, h))\) and \((m(f, g), h)\) belong to C and one has \(m(f, m(g, h)) = m(m(f, g), h)\).
\item
  For every \(f \in \mathrm{F}\), there exist elements \(o(f)\), \(t(f)\), and \(\overline{f} \in \mathrm{F}\), necessarily unique by virtue of (i), such that the pairs \((f, t(f))\), \((o(f), f)\), \((f, \overline{f})\), and \((\overline{f}, f)\) belong to C and one has \(m(f, t(f)) = f = m(o(f), f)\) and \(m(f, \overline{f}) = o(f)\).
\end{enumerate}

\begin{enumerate}
\def\labelenumi{\alph{enumi})}
\item
  Prove that, for every \(f \in F\) , we have \((\overline{f}, f) \in C\) and \(m(\overline{f}, f) = t(f)\) . Deduce that \((o(f), o(f)) \in C\) and \(m(o(f), o(f)) = o(f)\) . Also prove that \(o(\overline{f}) = t(f)\) and \(\overline{\overline{f}} = f\) .
\item
  Let S be the set of elements \(e\) of F such that \((e,e)\in\mathrm{C}\) and \(m(e,e)=e\); let us denote by \(o\) and \(t\) the maps from F to S given by \(f\mapsto o(f)\) and \(f\mapsto t(f)\).
\end{enumerate}

Prove that the quadruple \(\Gamma = (S, F, o, t)\) is a quiver and the map \(m\) is a composition law in \(\Gamma\) which makes it a groupoid. (This construction is due to H. BRANDT; cf. “Über eine Verallgemeinerung des Gruppenbegriffes”, Mathematische Annalen (1926), vol. 96, p. 360–366.)

\begin{enumerate}
\def\labelenumi{\arabic{enumi})}
\setcounter{enumi}{10}
\tightlist
\item
  Let G be a groupoid; let C be the set of composable pairs of arrows of G, and let \(m\colon\mathrm{C}\to\mathrm{Fl}(\mathrm{G})\) be the composition law of G.
\end{enumerate}

\begin{enumerate}
\def\labelenumi{\alph{enumi})}
\item
  Show that C and m satisfy hypotheses (i), (ii), (iii) of Exercise 10. Denote by Γ the groupoid provided by the construction in that exercise.
\item
  Prove that the pair consisting of the map \(e \mapsto o(e)\) from \(\operatorname{Som}(\Gamma)\) to \(\operatorname{Som}(\mathrm{G})\) and of the identity mapping of the set \(\operatorname{Fl}(\mathrm{G})\) into itself is an isomorphism of groupoids from \(\Gamma\) to G.
\end{enumerate}

This shows that a groupoid is determined by the set of its arrows and its composition law.

\begin{enumerate}
\def\labelenumi{\arabic{enumi})}
\setcounter{enumi}{11}
\tightlist
\item
  Let G be a group, let S be the set of subgroups of G, and let F be the set of subsets X of G such that there exist a subgroup H of G and an element g of G with X = gH.
\end{enumerate}

\begin{enumerate}
\def\labelenumi{\alph{enumi})}
\item
  Let \(\mu\) be the map from \(G \times G \times G\) to \(G\) defined by \(\mu(x,y,z)=xy^{-1}z\). Prove that a non-empty subset \(X\) of \(G\) belongs to the set F if and only if \(\mu(X\times X\times X)\subset X\).
\item
  Show that, by setting \(o(\mathrm{X}) = g^{-1}\mathrm{X}\) and \(t(\mathrm{X}) = \mathrm{X}g^{-1}\) for an arbitrary element g of X, one defines maps \(o: F \to S\) and \(t: F \to S\).
\item
  If \(X_{1}\) and \(X_{2}\) are elements of S such that \(t(\mathrm{X}_{1}) = o(\mathrm{X}_{2})\), we set \(m(\mathrm{X}_{1}, \mathrm{X}_{2}) = \mathrm{X}_{1}\mathrm{X}_{2}\). Show that the map m is a composition law in the quiver (S,F,o,t) which makes it a groupoid ("Baer groupoid"). It is denoted \(\mathrm{Ba}(\mathrm{G})\).
\item
  Show that the orbits of \(\mathrm{Ba}(\mathrm{G})\) are the conjugacy classes of subgroups of G.
\item
  Show that the isotropy group of Ba(G) at an element H of S is isomorphic to the quotient group N(H)/H.
\end{enumerate}

\setcounter{exercisegroup}{4}

\exercisegroup

\begin{enumerate}
\def\labelenumi{\arabic{enumi})}
\tightlist
\item
  Let G be a groupoid, X a set, and f a map from Som(G) to X.
\end{enumerate}

\begin{enumerate}
\def\labelenumi{\alph{enumi})}
\item
  Prove that there exists a pair \((K,\alpha)\), where K is a groupoid with vertex set X and \(\alpha: G \rightarrow K\) is a morphism of groupoids such that \(\mathrm{Som}(\alpha)=f\), satisfying the following universal property: for every groupoid H with vertex set X and every morphism \(\varphi:G\to H\) such that \(\mathrm{Som}(\varphi)=f\), there exists a unique morphism of groupoids \(\psi:K\to H\) such that \(\alpha\circ\psi=\varphi\). (If S is a set, we denote by \(S_{d}\) a groupoid with vertex set S whose orbits and isotropy groups are reduced to a single element. Let \(G'\) be the disjoint union groupoid of the family \((\mathrm{G},\mathrm{X}_{\mathrm{d}})\), let \(\varphi\) and \(\psi\) be the morphisms of the groupoid \(\mathrm{Som}(\mathrm{G})_{\mathrm{d}}\) to \(G'\) such that \(\mathrm{Som}(\varphi)\) is the canonical injection of \(\mathrm{Som}(\mathrm{G})\) into \(\mathrm{Som}(\mathrm{G}')\) and \(\mathrm{Som}(\psi)=f\), respectively. Set \(K=\mathrm{Coeg}(\varphi,\psi).\).)
\item
  Deduce from this a bijection of \(\operatorname{Hom}(G, f^{*}H)\) onto \(\operatorname{Hom}(K, H)\) for every groupoid H with vertex set X.
\end{enumerate}

\begin{enumerate}
\def\labelenumi{\arabic{enumi})}
\setcounter{enumi}{1}
\tightlist
\item
  Let H and G be groups, let \(\mathcal{H}\) and \(\mathcal{G}\) be the groupoids associated with them; let \(\varphi\) and \(\psi\) be morphisms of groupoids from \(\mathcal{H}\) to \(\mathcal{G}\).
\end{enumerate}

\begin{enumerate}
\def\labelenumi{\alph{enumi})}
\item
  Prove that the set of vertices of the cohomotoper (resp. of the coequalizer) of the pair \((\varphi,\psi)\) is reduced to a single element; we denote by K (resp. L) its isotropy group at this point.
\item
  Prove that the group K is the quotient of the free product \(G * \mathbf{Z}\) by the smallest normal subgroup of \(G * \mathbf{Z}\) that contains the elements \(\varphi(g)t\psi(g)^{-1}t^{-1}\), for \(g \in G\), where t denotes the element 1 of \(\mathbf{Z}\).
\item
  Prove that the group L is the quotient of the group G by the smallest normal subgroup containing the elements \(\varphi(g)\psi(g)^{-1}\), for \(g \in G\).
\end{enumerate}

\begin{enumerate}
\def\labelenumi{\arabic{enumi})}
\setcounter{enumi}{2}
\tightlist
\item
  Let I be a preordered set and let \((\mathrm{G}_{i})_{i\in\mathrm{I}}\) be a family of groupoids. For every pair \((i,j)\) of elements of I such that \(i\leqslant j\), let \(\varphi_{ji}\) be a groupoid morphism from \(G_{i}\) to \(G_{j}\). It is assumed that the \(\varphi_{ji}\) satisfy the following conditions (cf. E, III, p. 61):
\end{enumerate}

\begin{enumerate}
\def\labelenumi{(\roman{enumi})}
\item
  The relations \(i \leqslant j \leqslant k\) imply \(\varphi_{ki} = \varphi_{kj} \circ \varphi_{ji}\);
\item
  For every \(i \in \mathrm{I}\), \(f_{ii}\) is the identity morphism of \(\mathbf{G}_i\).
\end{enumerate}

Let G be the sum groupoid of the family \((\mathrm{G}_{i})\); let H be the sum groupoid of the family \((\mathrm{H}_{ji})_{i\leqslant j}\), where we have set \(H_{ji} = G_{i}\) for every pair \((i,j)\). Let \(\varphi\) (resp. \(\psi\)) be the morphism from H to G induced by the family \((\varphi_{ji})_{i\leqslant j}\) (resp. \((\mathrm{Id}_{\mathrm{G}_{i}})_{i\leqslant j}\)) and let C be the coequalizer of the pair \((\varphi,\psi)\). For \(i \in I\), let \(\varphi_{i}\) be the morphism from \(G_{i}\) to C, composed of the canonical morphism from \(G_{i}\) to G and the canonical morphism from G to C.

\begin{enumerate}
\def\labelenumi{\alph{enumi})}
\tightlist
\item
  Show that one has \(\varphi_j \circ \varphi_{ji} = \varphi_i\) for every pair \((i,j)\) such that \(i \leqslant j\). b) Show that, for every groupoid D and every family \((g_i)\), where \(g_i\) is a morphism \(\mathrm{G}_i\) to D such that \(g_j \circ \varphi_{ji} = g_i\) for every pair \((i,j)\) such that \(i \leqslant j\), there exists a unique morphism of groupoids \(g: \mathrm{C} \to \mathrm{D}\) such that \(g \circ \varphi_i = g_i\) for all \(i \in \mathrm{I}\).
\end{enumerate}

\chapter{Homotopy and Poincaré groupoid}

\section{HOMOTOPIES, HOMOTOPY EQUIVALENCES}

In this section and those that follow, I denotes the interval {[}0,1{]} of R.

\subsection{Homotopic continuous maps}

DEFINITION 1. --- Let X and Y be topological spaces and let f and g be continuous maps from X to Y. A homotopy connecting f to g is a continuous map \(\sigma: X \times I \to Y\) such that, for every \(x \in X\), we have \(\sigma(x,0) = f(x)\) and \(\sigma(x,1) = g(x)\). We say that f is homotopic to g if there exists a homotopy connecting f to g.

We say that \(f\) is the origin and \(g\) the term of the homotopy \(\sigma\) (cf. below, III, p. 257, remark 2).

Let A be a subset of X. We say that the homotopy \(\sigma\) is fixed on A if, for every \(a \in \mathrm{A}\), the map \(t \mapsto \sigma(a, t)\) from I to Y is constant. In this case, the origin and the term of \(\sigma\) coincide at every point of A.

Let X and Y be topological spaces. We say that homotopies \(\sigma\colon X\times I\to Y\) and \(\tau\colon X\times I\to Y\) are juxtaposable if the term of \(\sigma\) is the origin of \(\tau\), in other words if one has \(\sigma(x,1)=\tau(x,0)\) for all \(x\in X\). In this case, the map \(\sigma*\tau\) from \(X\times I\) to Y defined by

\[
(\sigma * \tau) (x, t) = \left\{ \begin{array}{l l} \sigma (x, 2 t) & \text {for} 0 \leqslant t \leqslant 1 / 2 \\ \tau (x, 2 t - 1) & \text {for} 1 / 2 \leqslant t \leqslant 1 \end{array} \right.\tag{1}
\]

is continuous (TG, I, p. 19, prop. 4) and is a homotopy connecting the origin of \(\sigma\) to the term of \(\tau\). It is called the juxtaposed homotopy of the homotopies \(\sigma\) and \(\tau\).

If \(\sigma\colon X\times\mathbf{I}\to Y\) is a homotopy, the map \(\overline{\sigma}\colon X\times\mathbf{I}\to Y\) defined by \((x,t)\mapsto\sigma(x,1-t)\) is a homotopy connecting the term of \(\sigma\) to the origin of \(\sigma\). We have \(\overline{\overline{\sigma}}=\sigma\). If \(\sigma\) and \(\tau\) are juxtaposable homotopies from \(X\times\mathbf{I}\) to Y, the homotopies \(\overline{\tau}\) and \(\overline{\sigma}\) are juxtaposable, and one has \(\overline{\sigma*\overline{\tau}}=\overline{\tau}*\overline{\sigma}\).

PROPOSITION 1. --- Let X and Y be topological spaces. The relation "f is homotopic to g" is an equivalence relation on the set \(\mathcal{C}(\mathrm{X};\mathrm{Y})\) of continuous maps from X to Y.

Let \(f\) be an element of \(\mathcal{C}(\mathrm{X};\mathrm{Y})\). The map \(f\circ\mathrm{pr}_{1}\colon\mathrm{X}\times\mathrm{I}\to\mathrm{Y}\) is a homotopy connecting \(f\) to \(f\); this relation is therefore reflexive.

Let \(f\) and \(g\) be elements of \(\mathcal{C}(\mathrm{X};\mathrm{Y})\) and \(\sigma\colon\mathrm{X}\times\mathbf{I}\to\mathrm{Y}\) a homotopy connecting \(f\) to \(g\). The map \(\overline{\sigma}\colon\mathrm{X}\times\mathbf{I}\to\mathrm{Y}\) is then a homotopy connecting \(g\) to \(f\); the relation under consideration is therefore symmetric.

Let us finally prove that it is transitive. If \(f\), \(g\), and \(h\) are elements of \(\mathcal{C}(\mathrm{X};\mathrm{Y})\), \(\sigma\) a homotopy connecting \(f\) to \(g\), and \(\tau\) a homotopy connecting \(g\) to \(h\), then \(\sigma\) and \(\tau\) are juxtaposable and \(\sigma*\tau\) is a homotopy connecting \(f\) to \(h\).

Let X and Y be topological spaces. The equivalence relation “f is homotopic to g” in \(\mathcal{C}(\mathrm{X};\mathrm{Y})\) (prop. 1) is called the homotopy relation. The quotient set of \(\mathcal{C}(\mathrm{X};\mathrm{Y})\) by this relation is denoted {[}X;Y{]}. Its elements are called homotopy classes of continuous maps from X to Y. The homotopy class of a continuous map \(f\colon \mathrm{X}\to \mathrm{Y}\) will often be denoted {[}f{]}.

PROPOSITION 2. — Let X, Y, and Z be topological spaces, f and \(f'\) be continuous maps from X to Y, and g and \(g'\) continuous maps from Y to Z. If f is homotopic to \(f'\) and if g is homotopic to \(g'\), then \(g \circ f\) is homotopic to \(g' \circ f'\).

Let \(\sigma\) be a homotopy connecting \(f\) to \(f'\) and \(\tau\) a homotopy connecting \(g\) to \(g'\). Then, the map \(\theta\colon X\times I\to Z\) defined by \(\theta(x,t)=\tau(\sigma(x,t),t)\) is a homotopy connecting \(g\circ f\) to \(g'\circ f'\).

Let X, Y, and Z be topological spaces. Given homotopy classes \(\varphi \in [\mathrm{X};\mathrm{Y}]\) and \(\psi \in [\mathrm{Y};\mathrm{Z}]\), the maps \(g \circ f: \mathrm{X} \to \mathrm{Z}\), where \(f \in \varphi\), \(g \in \psi\), all belong to the same homotopy class (prop. 2), which is denoted \(\psi \circ \varphi\) and is called the composite homotopy class of the classes \(\psi\) and \(\varphi\). The map {[}X;Y{]} \(\times\) {[}Y;Z{]} \(\to\) {[}X;Z{]} which associates \(\psi \circ \varphi\) to \((\varphi,\psi)\) is called the composition map.

Let \(\varphi\in[\mathrm{X};\mathrm{Y}]\); we have \(\varphi=\varphi\circ[\mathrm{Id}_{\mathrm{X}}]=[\mathrm{Id}_{\mathrm{Y}}]\circ\varphi\).

Let X, Y, Z, and T be topological spaces, and let \(\varphi \in [\mathrm{X};\mathrm{Y}]\), \(\psi \in [\mathrm{Y};\mathrm{Z}]\), and \(\chi \in [\mathrm{Z};\mathrm{T}]\). The homotopy class \(\chi \circ (\psi \circ \varphi)\) is equal to \((\chi \circ \psi) \circ \varphi\); it is denoted \(\chi \circ \psi \circ \varphi\).

PROPOSITION 3. --- Let X be a topological space and let \((\mathrm{Y}_{j})_{j\in\mathrm{J}}\) be a family of topological spaces. The map from \([X;\prod_{j\in\mathrm{J}}Y_{j}]\) to the product set \(\prod_{j\in\mathrm{J}}[X;Y_{j}]\) defined by \(\varphi\mapsto([\mathrm{pr}_{j}]\circ\varphi)_{j\in\mathrm{J}}\) is bijective.

Surjectivity follows immediately from I, p. 25, prop. 1. Let us prove injectivity. Let \(f\) and \(g\) be continuous maps from X to \(\prod_{j \in \mathrm{J}} \mathrm{Y}_j\). For all \(j \in \mathrm{J}\), set \(f_j = \mathrm{pr}_j \circ f\) and \(g_j = \mathrm{pr}_j \circ g\); suppose \(f_j\) is homotopic to \(g_j\) and let \(\sigma_j\) be a homotopy connecting \(f_j\) to \(g_j\). The map \(\sigma = (\sigma_j)\) from \(\mathrm{X} \times \mathbf{I}\) to \(\prod_{j \in \mathrm{J}} \mathrm{Y}_j\) is continuous (loc. cit.); this is a homotopy connecting \(f\) to \(g\), whence the proposition.

COROLLARY. --- Let \((\mathrm{X}_{j})_{j\in\mathrm{J}}\) and \((\mathrm{Y}_{j})_{j\in\mathrm{J}}\) be families of topological spaces with the same index set. For every \(j\in J\), let \(f_{j}\) and \(g_{j}\) be continuous maps from \(X_{j}\) to \(Y_{j}\). If, for every \(j\in J\), the maps \(f_{j}\) and \(g_{j}\) are homotopic, then so are the product maps \(f\colon(x_{j})\mapsto(f_{j}(x_{j}))\) and \(g\colon(x_{j})\mapsto(g_{j}(x_{j}))\) from \(\prod_{j\in J}X_{j}\) to \(\prod_{j\in J}Y_{j}\).

Suppose that we have \([f_j] = [g_j]\) for all \(j \in \mathbf{J}\). We have \([\mathrm{pr}_j] \circ [f] = [f_j \circ \mathrm{pr}_j]\) and \([\mathrm{pr}_j] \circ [g] = [g_j \circ \mathrm{pr}_j]\); therefore \([\mathrm{pr}_j] \circ [f] = [\mathrm{pr}_j] \circ [g]\) by prop. 2, whence \([f] = [g]\) by prop. 3.

\subsection{Pointed homotopies}

Let X and Y be topological spaces and let x be a point of X. A homotopy \(\sigma\colon X \times I \to Y\) is said to be pointed at x if it is fixed on \(\{x\}\), that is, if the map \(t \mapsto \sigma(x, t)\) from I to Y is constant. The juxtaposition of two homotopies pointed at x is pointed at x. The origin and the term of a homotopy pointed at x take the same value y at x; they are therefore pointed continuous maps from \((X, x)\) to \((Y, y)\).

Let \((X,x)\) and \((Y,y)\) be pointed topological spaces (I, p. 120, definition 1). The relation “f is connected to g by a homotopy pointed at x” is an equivalence relation on the set \(\mathcal{C}((X,x);(Y,y))\), called the pointed homotopy relation. The quotient set of \(\mathcal{C}((X,x);(Y,y))\) by this equivalence relation is denoted \([(X,x);(Y,y)]\). Its elements are called pointed homotopy classes of pointed continuous maps from \((X,x)\) to \((Y,y)\).

\pandocbounded{\includegraphics[keepaspectratio]{images/part002_b3045f29.jpg}}

The set \(\mathcal{C}((\mathrm{X},x);(\mathrm{Y},y))\) is a subset of \(\mathcal{C}(\mathrm{X};\mathrm{Y})\). Note that pointed homotopy is an equivalence relation finer than the relation induced by homotopy on \(\mathcal{C}(\mathrm{X};\mathrm{Y})\). It is in general a strictly finer relation (see III, p. 321, exerc. 1 and III, p. 234, example 3).

Let also \((Z,z)\) be a pointed topological space. Let f and \(f'\) be pointed continuous maps from \((X,x)\) to \((Y,y)\), and g and \(g'\) be pointed continuous maps from \((Y,y)\) to \((Z,z)\). If f and \(f'\) are related by a pointed homotopy \(\sigma\) at x, and if g and \(g'\) are related by a pointed homotopy \(\tau\) pointed at y, then \(g \circ f\) and \(g' \circ f'\) are connected by the homotopy \(\theta: X \times I \to Z\), \((u,t) \mapsto \tau(\sigma(u,t),t)\), which is pointed at x. As above, this allows one to define the composition map of \([(X,x);(Y,y)] \times [(Y,y);(Z,z)]\) to \([(X,x);(Z,z)]\). We leave it to the reader to formulate and prove for pointed maps and pointed homotopies the statements analogous to Prop. 3 and its corollary.

\subsection{Homotopy-equivalent spaces}

DEFINITION 2. --- Let X and Y be topological spaces. A homotopy class \(\varphi \in [\mathrm{X};\mathrm{Y}]\) is invertible if there exists a homotopy class \(\psi \in [\mathrm{Y};\mathrm{X}]\) such that \(\psi \circ \varphi = [\mathrm{Id}_{\mathrm{X}}]\) and \(\varphi \circ \psi = [\mathrm{Id}_{\mathrm{Y}}]\). We say that a continuous map is a homotopy equivalence if its homotopy class is invertible.

Let \(\varphi\in[X;Y]\) be an invertible homotopy class. There exists a unique homotopy class \(\psi\in[Y;X]\) having the properties of Definition 2. Indeed, let \(\psi,\psi'\) be classes having these properties; we have

\[
\psi = \psi \circ [ \mathrm{Id} _ {\mathrm{Y}} ] = \psi \circ \varphi \circ \psi^ {\prime} = [ \mathrm{Id} _ {\mathrm{X}} ] \circ \psi^ {\prime} = \psi^ {\prime}.
\]

This unique class is called the inverse of the homotopy class \(\varphi\) and is denoted \(\varphi^{-1}\).

Let Z be a topological space. We have \(\chi\circ\varphi\circ\varphi^{-1}=\chi\) for all \(\chi\in[Y;Z]\) and \(\theta\circ\varphi^{-1}\circ\varphi=\theta\) for all \(\theta\in[X;Z].\) It follows that the maps \(\chi\mapsto\chi\circ\varphi\) from {[}Y;Z{]} to {[}X;Z{]} and \(\theta\mapsto\theta\circ\varphi^{-1}\) from {[}X;Z{]} to {[}Y;Z{]} are inverse bijections of each other. Likewise, the map \(\chi\mapsto\varphi\circ\chi\) from {[}Z;X{]} to {[}Z;Y{]} is bijective and its inverse is the map \(\theta\mapsto\varphi^{-1}\circ\theta\) from {[}Z;Y{]} to {[}Z;X{]}.

Let X, Y, Z be topological spaces, and let \(\varphi \in [X;Y]\) and \(\psi \in [Y;Z]\) be invertible homotopy classes. Then the class \(\psi \circ \varphi\) is invertible, with inverse \(\varphi^{-1} \circ \psi^{-1}\). Indeed, we have

\[
\begin{array}{r l} & (\psi \circ \varphi) \circ (\varphi^ {- 1} \circ \psi^ {- 1}) = \psi \circ (\varphi \circ \varphi^ {- 1}) \circ \psi^ {- 1} \\ & \qquad = \psi \circ [ \mathrm{Id} _ {\mathrm{Y}} ] \circ \psi^ {- 1} \\ & \qquad = \psi \circ \psi^ {- 1} = [ \mathrm{Id} _ {\mathrm{Z}} ] \end{array}
\]

and, similarly, \((\varphi^{-1} \circ \psi^{-1}) \circ (\psi \circ \varphi) = [\mathrm{Id}_{\mathrm{X}}]\).

Let X and Y be topological spaces and let \(f: X \to Y\) be a homotopy equivalence. Every continuous map \(g: Y \to X\) whose class is the inverse of that of f is said to be a homotopy inverse (or inverse) of f up to homotopy. Such a map g is a homotopy equivalence.

A homeomorphism \(f\) is a homotopy equivalence, and \([f]^{-1} = [f^{-1}]\). The composite of two homotopy equivalences is a homotopy equivalence. The relation “X and Y are topological spaces and there exists a homotopy equivalence from X to Y” is an equivalence relation.

DEFINITION 3. --- Topological spaces X and Y are said to be homotopy equivalent if there exists a homotopy equivalence from X to Y.

Examples. --- 1) The empty topological space is homotopy equivalent only to itself.

\begin{enumerate}
\def\labelenumi{\arabic{enumi})}
\setcounter{enumi}{1}
\item
  Let X be a nonempty topological space. For X to be homotopy equivalent to a space reduced to a point, it is necessary and sufficient that there exist a constant map \(p: X \rightarrow X\) which is homotopic to the identity map of X. (It is then so for every constant map \(q \colon X \to X\) since \([p] = [p] \circ [q] = [\mathrm{Id}_{\mathrm{X}}] \circ [q] = [q]\).) Indeed, let \(P\) be a space reduced to a point, \(f \colon P \to X\) a map, and \(g \colon X \to P\) the unique map from X to P. We have \(g \circ f = \mathrm{Id}_{\mathrm{P}}\); for f to be a homotopy equivalence, it is necessary and sufficient that \(f \circ g\) be homotopic to \(\mathrm{Id}_{\mathrm{X}}\). Now, \(f \circ g\) is constant, with image \(f(P)\).
\item
  A pointed topological space (X,x) is said to be contractible, or the topological space X is said to be contractible at x, if there exists a pointed homotopy at x connecting \(Id_{X}\) to the constant map from X into X with image \(\{x\}\). Such a space is homotopy equivalent to a point. There exist, however, spaces homotopy equivalent to a point which are contractible at none of their points, and spaces contractible at one point, but not at every point (III, p. 321, exerc. 1).
\item
  Let E be the n-dimensional Euclidean space (or, more generally, a topological vector space over R) and let X be a subset of E. The set X is said to be star-shaped at a point x of X if, for every \(y \in X\) and every \(t \in I\), the point \(tx + (1 - t)y\) belongs to X. A convex subset (I, p. 122) of E is star-shaped at each of its points.
\end{enumerate}

A topological subspace X of E that is star-shaped at one of its points x is contractible at that point. Indeed, the map \(\sigma\colon X \times I \to X\) defined by \(\sigma(y,t) = tx + (1 - t)y\) is a pointed homotopy at x connecting \(Id_{X}\) to the constant map with image \(\{x\}\).

In particular, every interval of R, every convex subset of a Euclidean space or, more generally, of a topological vector space over R, is contractible at each of its points.

\begin{enumerate}
\def\labelenumi{\arabic{enumi})}
\setcounter{enumi}{4}
\tightlist
\item
  We will demonstrate later (TA, V) that Euclidean spheres \(S_n\), \(n\geqslant1\), are not homotopy equivalent to a point. The unit sphere of a separable infinite-dimensional Hilbert space is contractible at each of its points (EVT, V, p. 71, exerc. 13).
\end{enumerate}

\subsection{Relative homotopies}

Let X, Y be topological spaces, let A be a subspace of X, and let B be a subspace of Y.

Let \(f:\mathrm{X}\to\mathrm{Y}\) be a continuous map such that \(f(\mathrm{A})\subset\mathrm{B}\). We say that f is a homotopy equivalence of the pair (X,A) onto the pair (Y,B)

if there exists a continuous map \(g\colon Y\to X\) such that \(g(B)\subset A\), a homotopy \(\sigma\colon X\times I\to X\), fixed on A, relating \(\text{Id}_X\) to \(g\circ f\), and a homotopy \(\tau\colon Y\times I\to Y\), fixed on B, relating \(\text{Id}_Y\) to \(f\circ g\).

Suppose these conditions are satisfied. Then:

\begin{enumerate}
\def\labelenumi{\alph{enumi})}
\item
  The map f is a homotopy equivalence from X to Y, the map g is a homotopy equivalence from Y to X, and these homotopy equivalences are inverse to each other up to homotopy.
\item
  The map g is then a homotopy equivalence of the pair (Y,B) onto the pair (X,A), called an inverse of f up to homotopy.
\item
  The maps f and g induce, by restriction to subspaces, homeomorphisms from A onto B that are inverse to each other.
\end{enumerate}

We say that the pairs (X,A) and (Y,B) are homotopy equivalent if there exists a homotopy equivalence of the pair (X,A) onto the pair (Y,B). The relation “X, Y are topological spaces, A is a subspace of X, B is a subspace of Y, and the pairs (X,A) and (Y,B) are homotopy equivalent” is an equivalence relation.

\subsection{Retractions and contractions}

DEFINITION 4. --- Let X be a topological space and let A be a subset of X.

We call a retraction of X onto A a continuous map from X to A which is a retraction of the canonical injection of A into X. If there exists a retraction of X onto A, we say that X can be retracted onto A, or again that A is a retract of X.

Let X be a topological space and let A be a subspace of X.

Suppose there exists a retraction \(r\) of X onto A. The subspace A is identified with the set of points \(x \in \mathrm{X}\) such that \(x = r(x)\). If X is a separated space, then A is closed in X.

By the following lemma, for a subspace A to be a retract of X, it is necessary and sufficient that every continuous map from A to a topological space Y extends to a continuous map from X to Y.

Lemma 1. --- Let X and Y be topological spaces, and let \(f\colon X\to Y\) be a continuous map. The following conditions are equivalent:

\begin{enumerate}
\def\labelenumi{(\roman{enumi})}
\item
  For every topological space Z and every continuous map \(g: X \to Z\), there exists a continuous map \(g': Y \to Z\) such that \(g = g' \circ f\);
\item
  The map f is injective and has a continuous retraction;
\item
  The map f defines a homeomorphism from X onto its image \(f(\mathrm{X})\), which is a retract of Y.
\end{enumerate}

Suppose that \(f\) is injective, and let \(r: \mathrm{Y} \to \mathrm{X}\) be a continuous retraction of the map \(f\). For every continuous map \(g: \mathrm{X} \to \mathrm{Z}\), the map \(g' = g \circ r: \mathrm{Y} \to \mathrm{Z}\) satisfies \(g' \circ f = g \circ r \circ f = g\). This proves that (ii)⇒(i).

Suppose that condition (i) is satisfied. Let \(g \colon X \to X\) be the identity map. By hypothesis, there exists a continuous map \(g' \colon Y \to X\) such that \(g' \circ f = g\). This implies that the map \(f\) is injective and that \(g'\) is a continuous retraction of \(f\).

The equivalence of properties (ii) and (iii) is immediate.

DEFINITION 5. --- Let X be a topological space and let A be a subset of X. We call a contraction of X onto A a homotopy \(\sigma\colon X\times I\to X\) fixed on A whose origin is the identity map of X and whose term is a retraction of X onto A. If there exists a contraction of X onto A, we say that X can be contracted onto A, or again that A is a contraction of X.

Remark. --- In other words, a contraction of X onto A is a homotopy \(\sigma\colon X \times I \to X\) having the following properties:

\begin{enumerate}
\def\labelenumi{(\roman{enumi})}
\item
  \(\sigma(x,0)=x\) for all \(x\in X\);
\item
  \(\sigma(x,1)\in\mathrm{A}\) for all \(x\in\mathrm{X}\);
\item
  \(\sigma(x,t)=x\) for all \(x\in\mathbf{A}\) and every \(t\in\mathbf{I}\).
\end{enumerate}

For there to exist a contraction of X onto A, it is necessary and sufficient that the canonical injection of A into X be a homotopy equivalence of the pair (A,A) onto the pair (X,A).

Let \(a\) be a point of X. A contraction of X onto the subspace \(\{a\}\) is none other than a homotopy pointed at \(a\) connecting the map \(\mathrm{Id}_{\mathrm{X}}\) to the constant map with image \(\{a\}\). Consequently, X contracts onto \(\{a\}\) if and only if X is contractible at \(a\) (example 3 of III, p. 234).

Let X be a topological space and let A be a subset of X. Let \(\sigma\) be a contraction of X onto A. The map r from X to A defined by \(r(x) = \sigma(x,1)\) is a retraction of X onto A and \(\sigma\) is a homotopy connecting \(\mathrm{Id}_{\mathrm{X}}\) to \(r\). The relations \(r \circ i = \mathrm{Id}_{\mathrm{A}}\) and \(i \circ r = r\) then imply that the maps \(i\) and \(r\) are homotopy equivalences, inverse to each other up to homotopy.

DEFINITION 6. --- With the preceding notation, \(\sigma\) is said to be a strong contraction if, in addition, one has \(r(\sigma(x,t))=r(x)\) for all \(x\in X\) and all \(t\in I\).

Example. --- Let X be the complement of the origin in \(B_{n}\). The map \(X \times I\) to X given by \((x, t) \mapsto ((1 - t) + t \frac{1}{\|x\|})x\) is a strong contraction of X onto \(S_{n-1}\). The retraction of X onto \(S_{n-1}\) associated with it is the map given by \(x \mapsto x / \|x\|\).

Lemma 2. --- Let X be a topological space, let U be an open subset of X, and let \(\sigma\) be a strong contraction of X onto \(X-U\). Then \(\sigma(U\times I)\subset\overline{U}\). In particular, \(\sigma(U\times\{1\})\) is contained in the boundary of U.

Let \(x \in \mathrm{U}\). The set of real numbers \(t \in \mathbf{I}\) such that \(\sigma(x, t) \in \mathrm{U}\) is open in \(\mathbf{I}\) and contains 0; let \(s\) be its upper bound. We have \(\sigma(x, s) \in \overline{\mathrm{U}}\). If \(s < 1\), \(\sigma(x, s) \notin \mathrm{U}\) by definition of \(s\); the same holds if \(s = 1\), since \(\sigma(x, 1) \in \mathrm{X} - \mathrm{U}\). Consequently, \(\sigma(x, s) \in \mathrm{Fr}(\mathrm{U})\). By definition of a strong contraction, we then have \(\sigma(x, s) = \sigma(\sigma(x, s), 1) = \sigma(x, 1)\); in particular, \(\sigma(x, 1) \in \overline{\mathrm{U}}\). Consequently, \(\sigma(x, 1)\) belongs to the boundary \(\overline{\mathrm{U}} \cap (\mathrm{X} - \mathrm{U})\) of \(\mathrm{U}\).

Let \(t \in \mathbf{I}\); if \(\sigma(x,t) \notin \mathrm{U}\), we also get \(\sigma(x,t) = \sigma(\sigma(x,t),1) = \sigma(x,1)\), therefore \(\sigma(x,t) \in \overline{\mathrm{U}}\), whence the lemma.

\subsection{Cylinder of a map}

Let X and Y be topological spaces, and let f be a continuous map from X to Y. Denote by U the topological sum space of the space \(U_{1} = X \times I\) and the space \(U_{2} = Y\), and identify \(U_{1}\) and \(U_{2}\) with subspaces of U via the canonical injections. Let R be the finest equivalence relation on U for which the points \((x,1)\) of \(U_{1}\) and \(f(x)\) of \(U_{2}\) are equivalent, for every \(x \in X\).

DEFINITION 7. --- We call the cylinder of the map f the quotient topological space U/R and denote it by \(\mathrm{Cyl}(f)\).

Denote by \(\alpha_{f}\colon\mathrm{X}\times\mathbf{I}\to\mathrm{Cyl}(f)\) and \(\beta_{f}\colon\mathrm{Y}\to\mathrm{Cyl}(f)\) the restrictions to \(\mathrm{U}_{1}\) and \(\mathrm{U}_{2}\) of the canonical surjection from U onto U/R. The map \(\alpha_{f}\) is a homotopy and its term is \(\beta_{f}\circ f\).

PROPOSITION 4 (Universal property of cylinders)

Let Z be a topological space, let \(\beta: \mathrm{Y} \to \mathrm{Z}\) be a continuous map and let \(\alpha: \mathrm{X} \times \mathbf{I} \to \mathrm{Z}\) be a homotopy whose term is \(\beta \circ f\). There exists a unique continuous map \(\varphi\) from \(\operatorname{Cyl}(f)\) to Z such that \(\alpha = \varphi \circ \alpha_f\) and \(\beta = \varphi \circ \beta_f\).

The map \(\psi\) from U to Z which coincides with \(\alpha\) on \(\times\) and with \(\beta\) on Y is continuous. We have \(\alpha(x,1)=\beta(f(x))\) for all \(x\in X\), therefore \(\psi\) passes to the quotient to define a continuous map \(\varphi\) from \(\mathrm{Cyl}(f)\) to Z such that \(\alpha=\varphi\circ\alpha_{f}\) and \(\beta=\varphi\circ\beta_{f}\). Since the images of \(\alpha_{f}\) and \(\beta_{f}\) cover \(\mathrm{Cyl}(f)\), \(\varphi\) is the only map from \(\mathrm{Cyl}(f)\) to Z satisfying these relations.

Example 1. --- Let X' and Y' be topological spaces and \(f': \mathrm{X}' \to \mathrm{Y}'\) a continuous map. Let \(u: \mathrm{X} \to \mathrm{X}'\) and \(v: \mathrm{Y} \to \mathrm{Y}'\) be continuous maps such that \(f' \circ u = v \circ f\). There exists a unique continuous map \(w: \operatorname{Cyl}(f) \to \operatorname{Cyl}(f')\) such that \(w \circ \alpha_f = \alpha_{f'} \circ (u \times \operatorname{Id}_\mathbf{I})\) and \(w \circ \beta_f = \beta_{f'} \circ v\).

By prop. 4, applied to Z=Y and \(\beta = \operatorname{Id}_{Y}\), there exists a unique continuous map \(\gamma_{f}\colon\operatorname{Cyl}(f)\to\operatorname{Y}\) such that \(\gamma_{f}(\alpha_{f}(x,s))=f(x)\) and \(\gamma_{f}(\beta_{f}(y))=y\) for \(x\in X\), \(s\in I\), and \(y\in Y\).

PROPOSITION 5. --- The map \(\alpha_{f}\) induces a homeomorphism of X × {[}0,1{[} onto an open subset of \(\mathrm{Cyl}(f)\). The map \(\beta_{f}\) defines a homeomorphism of Y onto the complement of this open set. The map \(\beta_{f} \circ \gamma_{f}\) is a continuous retraction of \(\mathrm{Cyl}(f)\) onto \(\beta_{f}(\mathrm{Y})\).

The set X × {[}0,1{[} is an open subset of U, saturated for the relation R, and the equivalence relation induced by R in X × {[}0,1{[} is the equality relation. The first assertion follows from this (TG, I, p. 23, cor. 10).

The complement of \(\alpha_{f}(\mathrm{X}\times[0,1])\) is \(\beta_{f}(\mathrm{Y})\). Since one has \(\gamma_{f}\circ\beta_{f}=\mathrm{Id}_{\mathrm{Y}}\), \(\beta_{f}\) defines a homeomorphism of Y onto \(\beta_{f}(\mathrm{Y})\) and \(\beta_{f}\circ\gamma_{f}\) is a continuous retraction of the canonical injection of \(\beta_{f}(\mathrm{Y})\) into \(\mathrm{Cyl}(f)\).

The closed subspace \(\beta_{f}(\mathrm{Y})\) of \(\operatorname{Cyl}(f)\) is called the base of the cylinder of \(f\). The map \(\beta_{f} \circ \gamma_{f}\) is called the canonical retraction of \(\operatorname{Cyl}(f)\) onto its base.

Consider the maps \(\sigma_{1}\colon\mathrm{U}_{1}\times\mathbf{I}\to\mathrm{Cyl}(f)\) and \(\sigma_{2}\colon\mathrm{U}_{2}\times\mathbf{I}\to\mathrm{Cyl}(f)\) defined by

\[
\begin{array}{c c} \sigma_ {1} ((x, s), t) = \alpha_ {f} (x, (1 - t) s + t) & \text {for} (x, s) \in \mathrm{X} \times \mathbf {I} \text {and} t \in \mathbf {I} \\ \sigma_ {2} (y, t) = \beta_ {f} (y) & \text {for} y \in \mathrm{Y} \text {and} t \in \mathbf {I}. \end{array}
\]

They are continuous. For \(x \in X\) and \(t \in I\) , we have

\[
\sigma_ {1} ((x, 1), t) = \alpha_ {f} (x, 1) = \beta_ {f} (f (x)) = \sigma_ {2} (f (x), t).
\]

There therefore exists a unique map \(\sigma_{f}\) from \(\mathrm{Cyl}(f)\times\mathbf{I}\) to \(\mathrm{Cyl}(f)\) such that \(\sigma_{f}\circ(\alpha_{f}\times\mathrm{Id}_{\mathbf{I}})=\sigma_{1}\) and \(\sigma_{f}\circ(\beta_{f}\times\mathrm{Id}_{\mathbf{I}})=\sigma_{2}\). The map \(\sigma_{f}\) is continuous (I, p. 19, prop. 10).

Proposition 6. — The map \(\sigma_{f}\colon\mathrm{Cyl}(f)\times\mathbf{I}\to\mathrm{Cyl}(f)\) is a strong contraction of \(\mathrm{Cyl}(f)\) onto its base. Its term is the canonical retraction of \(\mathrm{Cyl}(f)\) onto its base.

The map \(\sigma_f\) is a homotopy. The relations \(\sigma_f(\alpha_f(x,s),0) = \alpha_f(x,s)\) and \(\sigma_f(\beta_f(y),0) = \beta_f(y)\) imply that the origin of \(\sigma_f\) is the identity map of \(\mathrm{Cyl}(f)\). Denote by \(r_f\) the term of \(\sigma_f\). The relations

\[
\sigma_ {f} \left(\alpha_ {f} (x, s), 1\right) = \alpha_ {f} (x, 1) = \beta_ {f} (f (x)) = \left(\beta_ {f} \circ \gamma_ {f}\right) \left(\alpha_ {f} (x, s)\right)
\]

and \(\sigma_{f}(\beta_{f}(y),1) = \beta_{f}(y) = (\beta_{f}\circ \gamma_{f})(\beta_{f}(y))\) imply that \(r_f\) is the canonical retraction \(\beta_{f}\circ \gamma_{f}\) of \(\mathrm{Cyl}(f)\) onto its base.

For \((x,s)\in \mathbf{X}\times \mathbf{I}\) and \(t\in \mathbf{I}\) , we have

\[
r _ {f} \left(\sigma_ {f} \left(\alpha_ {f} (x, s), t\right)\right) = r _ {f} \left(\alpha_ {f} (x, s (1 - t) + t)\right) = \beta_ {f} (f (x)) = r _ {f} \left(\alpha_ {f} (x, s)\right).
\]

For \(y \in Y\) and \(t \in I\), we have \(r_{f}(\sigma_{f}(\beta_{f}(y), t)) = r_{f}(\beta_{f}(y))\). Consequently, \(\sigma_{f}\) is a strong contraction of \(\operatorname{Cyl}(f)\) onto its base.

The map \(\sigma_{f}\) is called the canonical contraction of the cylinder of \(f\) onto its base.

Remarks. --- 1) Let A be a subset of X × I and let A₁ be the set of points x ∈ X such that (x,1) ∈ A. Then we have \(\bar{\alpha}_{f}^{-1}(\alpha_{f}(A)) = A \cup f^{-1}(f(A_{1})) \times \{1\}\) and \(\beta_{f}^{-1}(\alpha_{f}(A)) = f(A_{1})\). Consequently, the map \(\alpha_{f}\) is closed (resp. open) if f is.

\begin{enumerate}
\def\labelenumi{\arabic{enumi})}
\setcounter{enumi}{1}
\tightlist
\item
  Suppose that the map \(f\) is proper. Let \(P\) be a point of \(\operatorname{Cyl}(f)\). If \(P = \alpha_f(x,t)\), with \(0 \leqslant t < 1\) and \(x \in X\), we have \(\overline{\alpha}_f^1(P) = (x,t)\). In the contrary case, there exists \(y \in Y\) such that \(P = \beta_f(y)\)
\end{enumerate}

and \(\bar{\alpha}_{f}^{1}(P) = \bar{f}^{1}(y) \times \{1\}\). This proves that the fibers of the map \(\alpha_{f}\) are quasi-compact. The map \(\alpha_{f}\) is then proper (TG, I, p. 75, th. 1), since it is closed.

By prop. 5 and TG, I, p. 72, prop. 2, the map \(\beta_f\) is itself proper. Consequently, if f is proper, the canonical surjection from U onto \(\mathrm{Cyl}(f)\) is proper, hence, in particular, universally strict (I, p. 20, corollary).

\begin{enumerate}
\def\labelenumi{\arabic{enumi})}
\setcounter{enumi}{2}
\tightlist
\item
  If the spaces X and Y are separated, then so is the cylinder of f. Indeed, let z and \(z'\) be two distinct points of \(\mathrm{Cyl}(f)\) and let us prove that they possess disjoint neighborhoods. We distinguish three cases.
\end{enumerate}

- There exists \((x,t)\) and \((x',t') \in \mathrm{X} \times [0,1[\) such that \(z = \alpha_f(x,t)\) and \(z' = \alpha_f(x',t')\) .

In this case, the assertion follows from the fact that the space \(X \times [0,1]\) is separated (TG, I, p. 54, prop. 7) and that the map \(\alpha_{f}\) induces a homeomorphism from this space onto an open subspace of \(\operatorname{Cyl}(f)\).

- There exists \((x,t)\in \mathrm{X}\times [0,1[\) such that \(z = \alpha_{f}(x,t)\) and \(y^{\prime}\in \mathrm{Y}\) such that \(z^{\prime} = \beta_{f}(y^{\prime})\).

Then, \(\alpha_{f}(\mathrm{X}\times [0,\frac{t + 1}{2}])\) and \(\operatorname {Cyl}(f)\mathbf{-}\alpha_{f}(\mathrm{X}\times [0,\frac{t + 1}{2}])\) are disjoint open neighborhoods of \(z\) and \(z^{\prime}\) in \(\operatorname {Cyl}(f)\).

- There exist \(y\) and \(y'\in Y\) such that \(z=\beta_f(y)\) and \(z'=\beta_f(y')\).

In this case, \(y\neq y'\); since Y is separated, there exists an open neighborhood V of y in Y and an open neighborhood \(V'\) of \(y'\) in Y such that \(V\cap V'=\varnothing\). Then, \((\beta_f\circ\gamma_f)^{-1}(V)\) and \((\beta_f\circ\gamma_f)^{-1}(V')\) are disjoint open subsets of \(\operatorname{Cyl}(f)\) containing respectively \(\beta_f(y)\) and \(\beta_f(y')\).

\subsection{The homotopy extension property}

DEFINITION 8. --- Let X be a topological space and let A be a subset of X. The pair (X,A) is said to have the homotopy extension property if, for every topological space Y, every continuous map \(f:\mathrm{X}\to\mathrm{Y}\), and every homotopy \(\sigma:\mathrm{A}\times\mathbf{I}\to\mathrm{Y}\) whose term is the map \(f|_{\mathrm{A}}\), there exists a homotopy \(\tau:\mathrm{X}\times\mathbf{I}\to\mathrm{Y}\) extending \(\sigma\) and whose term is the map f.

Remark 1. --- Let X be a topological space and let A be a subset of X such that the pair (X,A) has the homotopy extension property. Let Y be a topological space, let \(f:\mathrm{X}\to\mathrm{Y}\) be a continuous map, and let \(\sigma:\mathrm{A}\times\mathrm{I}\to\mathrm{Y}\) be a homotopy whose origin is the map \(f|_{\mathrm{A}}\). The map \(\overline{\sigma}:\mathrm{A}\times\mathrm{I}\to\mathrm{Y}\) defined by \((a,t)\mapsto\sigma(a,1-t)\) is a homotopy with term \(f|_{\mathrm{A}}\); let \(\tau:\mathrm{X}\times\mathrm{I}\to\mathrm{Y}\) be a homotopy with term f which extends \(\overline{\sigma}\). The map \(\overline{\tau}:\mathrm{X}\times\mathrm{I}\to\mathrm{Y}\) given by \((x,t)\mapsto\tau(x,1-t)\) is then a homotopy extending \(\sigma\) and whose origin is the map f.

Let X be a topological space, let A be a subspace of X, and let \(i: A \to X\) be the canonical injection. Denote by \(\alpha_{i}: A \times I \to \text{Cyl}(i)\) and \(\beta_{i}: X \to \text{Cyl}(i)\) the canonical maps. Let \(j: \text{Cyl}(i) \to X \times I\) be the unique continuous map such that \(j(\alpha_{i}(a, s)) = (i(a), s)\) and \(j(\beta_{i}(x)) = (x, 1)\) for \(a \in A\), \(s \in I\), and \(x \in X\). It is injective; its image is the subspace \((\mathbf{A} \times \mathbf{I}) \cup (\mathbf{X} \times \{1\})\) of \(X \times I\). The map j is closed if A is closed in X. It is not always strict (III, p. 325, exerc. 17).

PROPOSITION 7. --- With the above notations, the following assertions are equivalent:

\begin{enumerate}
\def\labelenumi{(\roman{enumi})}
\item
  The pair (X, A) has the homotopy extension property;
\item
  For every topological space Y and every continuous map \(g\colon \operatorname{Cyl}(i) \to \operatorname{Y}\), there exists a continuous map \(g'\colon X \times I \to Y\) such that \(g = g' \circ j\);
\item
  The injection j has a continuous retraction;
\item
  The map j is strict and there exists a contraction of X × I onto the image of j.
\end{enumerate}

Suppose that the pair (X,A) has the homotopy extension property. Let \(g:\operatorname{Cyl}(i)\to Y\) be a continuous map; set \(\sigma=g\circ\alpha_i\) and \(f=g\circ\beta_i\). The map \(f:X\to Y\) is continuous and \(\sigma:A\times I\to Y\) is a homotopy whose term is the map \(f|A\). There therefore exists a homotopy \(g':X\times I\to Y\) with term f which extends \(\sigma\). We have \(g'\circ j(\alpha_i(a,s))=g'(a,s)=\sigma(a,s)=g(\alpha_i(a,s))\) and \(g'\circ j(\beta_i(x))=g'(x,1)=f(x)=g(\beta_i(x))\) for \(a\in A\), \(s\in I\), and \(x\in X\). Consequently, \(g'\circ j=g\), whence (ii).

Conversely, suppose that assertion (ii) holds and let us prove that the pair (X,A) has the homotopy extension property. Let Y be a topological space, let \(f:X\to Y\) be a continuous map and let \(\sigma:A\times I\to Y\) be a homotopy whose term is equal to \(f|A\). There exists a unique continuous map \(g:\mathrm{Cyl}(i)\to\mathrm{Y}\) such that \(g\circ\alpha_i=\sigma\) and \(g\circ\beta_i=f\) (III, p. 238, prop. 4). Every map \(g':\mathrm{X}\times\mathbf{I}\to\mathrm{Y}\) such that \(g=g'\circ j\) is then a homotopy with term f which extends \(\sigma\).

  The equivalence of assertions (ii) and (iii) is a special case of lemma 1 of III, p. 235.

Suppose that the injection \(j\colon \mathrm{Cyl}(i)\to \mathrm{X}\times \mathbf{I}\) admits a continuous retraction \(r\). The map \(j\) defines a homeomorphism from \(\mathrm{Cyl}(i)\) onto the subspace \(\mathrm{T} = (\mathrm{A}\times \mathbf{I})\cup (\mathrm{X}\times \{1\})\) of \(\mathrm{X}\times \mathbf{I}\) (loc. cit.). It is therefore strict.

Let us set \(\rho = \mathrm{pr}_1\circ j\circ r\) and \(\theta = \mathrm{pr}_2\circ j\circ r\). If \(x\in \mathrm{A}\) or if \(s = 1\), we have \(\rho (x,s) = x\) and \(\theta (x,s) = s\). For \((x,s)\in \mathrm{X}\times \mathbf{I}\) and \(t\in \mathbf{I}\), set

\[
\sigma ((x, s), t) = (\rho (x, (1 - t) + s t), (1 - t) s + t \theta (x, s)).
\]

The map \(\sigma\colon(\mathrm{X}\times\mathbf{I})\times\mathbf{I}\to\mathrm{X}\times\mathbf{I}\) is continuous. For \((x,s)\in\mathrm{X}\times\mathbf{I}\) , we have

\[
\sigma ((x, s), 0) = (\rho (x, 1), s) = (x, s)
\]

and

\[
\sigma ((x, s), 1) = (\rho (x, s), \theta (x, s)) = j \circ r (x, s);
\]

For \(x\in \mathbf{X}\) and \(t\in \mathbf{I}\), we have

\[
\sigma ((x, 1), t) = (\rho (x, 1), (1 - t) + t \theta (x, 1)) = (x, 1),
\]

while, for \((x,s)\in\mathbf{A}\times\mathbf{I}\) and \(t\in\mathbf{I}\) , we have

\[
\sigma ((x, s), t) = (x, (1 - t) s + t s) = (x, s).
\]

Therefore, \(\sigma\) is a contraction of \(X\times I\) onto T. This shows that assertion (iii) implies assertion (iv).

The implication (iv)⇒(iii) is evident.

Remark 2. --- Let X be a separated topological space and let A be a subset of X such that the pair (X,A) has the homotopy extension property. Denote by \(i\colon A\to X\) the canonical injection. Let \(j\colon\operatorname{Cyl}(i)\to X\times\mathbf{I}\) be the canonical injection and let \(r\) be a continuous retraction of the map j. The subspace \(j(\operatorname{Cyl}(i))\) is equal to the set of pairs \((x,t)\in X\times\mathbf{I}\) such that \(j(r(x,t))=(x,t)\). Since the space \(X\times\mathbf{I}\) is separated, the subset \(j(\operatorname{Cyl}(i))\) is closed in \(X\times\mathbf{I}\). The set A, equal to the set of \(x\in X\) such that \((x,0)\in j(\operatorname{Cyl}(i))\), is then a closed subset of X.

Lemma 3. --- Let X and Y be topological spaces, let \(p \colon X \to Y\) be a proper and open continuous map, and let \(f \colon X \to \mathbf{R}\) be a continuous map. The map \(g \colon Y \to \overline{\mathbf{R}}\) given by \(y \mapsto \sup_{x\in p^{-1}(y)} f(x)\) is continuous.

Let \(b \in \mathbf{Y}\). Since \(p\) is proper, its fiber \(\overline{p}^1(b)\) is a quasi-compact space; we therefore have \(g(b) \in \mathbf{R} \cup \{-\infty\}\).

Let \(m \in R\) be such that \(g(b) < m\). For all \(a \in \overline{p}^{1}(b)\), we have \(f(a) \leqslant g(b) < m\); let \(V_{a}\) be a neighborhood of a in X such that \(f(x) < m\) for all \(x \in V_{a}\). The union V of the sets \(V_{a}\) is a neighborhood of \(\overline{p}^{1}(b)\) in X. By lemma 5 (I, p. 75), there exists a neighborhood W of b in Y such that \(\overline{p}^{1}(W) \subset V\). For all \(y \in W\), we have \(g(y) \leqslant m\). This proves that g is upper semicontinuous at b.

Let us now prove that \(g\) is lower semicontinuous at \(b\). We may assume \(g(b)\in\mathbf{R}\). Let \(m\in\mathbf{R}\) be such that \(m<g(b)\). Let \(a\in\overline{p}^1(b)\) such that \(m<f(b)\); let then V be a neighborhood of a in X such that \(f(x)>m\) for all \(x\in V\). It follows that \(g(y)>m\) for all \(y\in p(V)\). Since p is open, \(p(V)\) is a neighborhood of b, so that g is lower semicontinuous.

The lemma is thus proved.

THEOREM 1. --- Let X be a topological space and let A be a closed subspace of X; denote by i : A → X the canonical injection. The following assertions are equivalent:

\begin{enumerate}
\def\labelenumi{(\roman{enumi})}
\item
  The pair (X, A) has the homotopy extension property;
\item
  There exists a continuous map \(\varphi\colon\mathrm{X}\to\mathbf{I}\) such that \(\mathrm{A}=\overline{\varphi}^{1}(0)\) and a homotopy \(\sigma\colon\mathrm{X}\times\mathbf{I}\to\mathrm{X}\) fixed on A whose term is the identity map of X and such that \(\sigma(x,0)\in\mathrm{A}\) for every point \(x\in\mathrm{X}\) such that \(\varphi(x)\neq1\).
\item
  There exists a continuous map \(\varphi\colon\mathbf{X}\to\mathbf{R}_{+}\) such that \(\mathrm{A}=\overline{\varphi}^{1}(0)\) and a homotopy \(\sigma\colon\overline{\varphi}^{1}(\mathbf{I})\times\mathbf{I}\to\mathrm{X}\), fixed on A, such that \(\sigma(x,1)=x\) and \(\sigma(x,0)\in\mathrm{A}\) for all \(x\in\overline{\varphi}^{1}(\mathbf{I})\).
\item
  There exists a continuous map \(\varphi\colon X\to\mathbf{R}_{+}\) such that \(A\subset\overline{\varphi}^{1}(0)\) and a continuous map
\end{enumerate}

\[
\sigma \colon \{(x, t) \in \mathrm{X} \times \mathbf {I} \mid t + \varphi (x) \geqslant 1 \} \to \mathrm{X}
\]

such that \(\sigma(x,1) = x\) for all \(x \in X\) and \(\sigma(x,1 - \varphi(x)) \in A\) for all \(x \in X\) such that \(\varphi(x) \leqslant 1\).

We denote by i the canonical injection of A into X, \(\alpha_{i}\colon A\times I\to\operatorname{Cyl}(i)\) and \(\beta_{i}\colon X\to\operatorname{Cyl}(i)\) the canonical maps, and by j the map from \(\operatorname{Cyl}(i)\) to \(X\times I\) such that \(j(\alpha_{i}(x,s))=(x,s)\) for \((x,s)\in A\times I\) and \(j(\beta_{i}(x))=(x,1)\) for \(x\in X\).

Suppose that the pair (X,A) has the homotopy extension property and let \(r:X\times I\to\operatorname{Cyl}(i)\) be a continuous retraction of the map j (III, p. 241, prop. 7). Denote by \(\sigma\) the continuous map \(\operatorname{pr}_1\circ j\circ r\) from \(X\times I\) to X. For every \((x,t)\in X\times I\), we have \(|\operatorname{pr}_2(j(r(x,t)))-t|\leqslant1\). Since I is compact, the first projection \(\operatorname{pr}_1:X\times I\to X\) is a proper map; it is also open. By lemma 3, the map \(\varphi:X\to I\) given by

\[
x \mapsto \sup _ {t \in \mathbf {I}} \left(\left| \operatorname{pr} _ {2} (j (r (x, t))) - t \right|\right)
\]

is therefore continuous.

For \(x \in \mathrm{X}\), we have \(\sigma(x,1) = x\). Let \(x \in \mathrm{X}\) such that \(\sigma(x,0) \notin \mathrm{A}\); then \(\mathrm{pr}_2(j(r(x,0))) = 1\) and \(\varphi(x) \geqslant 1\). For \(x \in \mathrm{A}\) and \(t \in \mathbf{I}\), we have \(j(r(x,t)) = (x,t)\); therefore, \(\sigma(x,t) = x\) and \(\varphi(x) = 0\). Conversely, let \(x \in \mathrm{X}\) such that \(\varphi(x) = 0\). We therefore have \(\mathrm{pr}_2(j(r(x,t)) \leqslant t\) for all \(t \in \mathbf{I}\); if \(t < 1\), this implies \(j(r(x,t)) \in \mathrm{A} \times \mathbf{I}\); since A is closed in X, we therefore have \(j(r(x,1)) \in \mathrm{A} \times \mathbf{I}\), and therefore \(x \in \mathrm{A}\).

This shows that (i) implies (ii).

Let \(\varphi\) and \(\sigma\) be maps satisfying the properties of assertion (ii). Set \(\varphi_1 = 2\varphi\) and let \(\sigma_1\) be the restriction of \(\sigma\) to \(\overline{\varphi}_1^1 (\mathbf{I})\times \mathbf{I}\). We have \(\overline{\varphi}_1^1 (0) = \mathrm{A}\); for \(a\in \mathrm{A}\), \(\sigma_1(a,t) = a\) for all \(a\in \mathrm{A}\). Let \(x\in \mathrm{X}\) such that \(\varphi_{1}(x)\leqslant 1\); we have \(\sigma_1(x,1) = x\); moreover, \(\varphi (x) = \varphi_{1}(x) / 2 < 1\), therefore \(\sigma_{1}(x,0)\in \mathrm{A}\). Thus, (ii) implies (iii).

Let us prove that (iii) implies (iv). Let \(\varphi\colon X\to\mathbf{R}\) and \(\sigma\colon\overline{\varphi}^{1}(\mathbf{I})\times\mathbf{I}\to X\) as in the statement; set \(B=\overline{\varphi}^{1}(\mathbf{I})\) and \(C=\overline{\varphi}^{1}([1,+\infty[)\) .

Let \(u_1\) be the map from B × I to X such that \(u_1(x,t) = \sigma(x,1 - (1 - t)/2\varphi(x))\) if \(t + 2\varphi(x) \geqslant 1\) and \(\varphi(x) > 0\), and \(u_1(x,t) = \sigma(x,0)\) otherwise. By Lemma 4 below, it is continuous.

Let \(u_2\) be the map from B × I to X such that \(u_2(x,t)=\sigma(x,\sup(0,1-2(1-t)(1-\varphi(x))))\). It is continuous. For \(x\in\mathrm{B}\) such that \(\varphi(x)=1/2\) and \(t\in\mathbf{I}\), we have \(u_1(x,t)=\sigma(x,t)=u_2(x,t)\). Denote by \(u:\mathrm{B}\times\mathbf{I}\to\mathrm{X}\) the map such that \(u(x,t)=u_1(x,t)\) if \(\varphi(x)\leqslant1/2\) and \(u(x,t)=u_2(x,t)\) if \(1/2<x\leqslant1\); it is continuous because its restrictions to the closed subspaces \(\overline{\varphi}^{1}([0,1/2])\times\mathbf{I}\) and \(\overline{\varphi}^{1}([1/2,1])\times\mathbf{I}\) of B×I are continuous (TG, I, p. 19, prop. 4).

For \(x\in X\) such that \(\varphi(x)=1\) and \(t\in\mathbf{I}\), we have \(u(x,t)=u_2(x,t)=\sigma(x,1)=x\). There therefore exists a unique map \(\tau:X\times I\to X\) which coincides with u on \(B\times I\) and with the map \(\mathrm{pr}_1\) on \(C\times I\); it is continuous (loc. cit.).

Let \(x \in X\); one verifies that \(\tau(x,1) = 1\). If, moreover, \(x \in A\), then \(\varphi(x) = 0\), therefore \(\tau(x,t) = u_{1}(x,t) = \sigma(x,0) = x\). Finally, if \(2\varphi(x) \leqslant 1\), then \(\varphi(x) \leqslant 1/2\), therefore \(\tau(x,1 - 2\varphi(x)) = \sigma(x,0) \in A\).

This proves that assertion (iv) holds.

Finally, suppose assertion (iv) is satisfied and let us prove that the pair (X, A) has the homotopy extension property.

Denote by \(C_{1}\) (resp. \(C_{2}\)) the set of pairs \((x,t)\in\mathbf{X}\times\mathbf{I}\) such that \(t+\varphi(x)\leqslant1\) (resp. \(t+\varphi(x)\geqslant1\)). These are closed sets. For \((x,t)\in\mathrm{C}_{1}\), we have \(\sigma(x,1-\varphi(x))\in\mathrm{A}\); let \(\rho_{1}\colon\mathrm{C}_{1}\to\mathrm{Cyl}(i)\) be the map given by \((x,t)\mapsto\alpha_{i}(\sigma(x,1-\varphi(x)),t+\varphi(x))\); it is continuous. Let also \(\rho_{2}\colon\mathrm{C}_{2}\to\mathrm{Cyl}(i)\) be the continuous map given by \((x,t)\mapsto\beta_{i}(\sigma(x,t))\). For \((x,t)\in\mathrm{C}_{1}\cap\mathrm{C}_{2}\), we have \(t+\varphi(x)=1\), therefore

\[
\rho_ {1} (x, t) = \alpha_ {i} (\sigma (x, 1 - \varphi (x)), 1) = \beta_ {i} (\sigma (x, 1 - \varphi (x)) = \rho_ {2} (x, t).
\]

There therefore exists a unique map \(\rho\colon X\times\mathbf{I}\to\mathrm{Cyl}(i)\) which coincides with \(\rho_{1}\) on \(C_{1}\) and with \(\rho_{2}\) on \(C_{2}\); it is continuous (TG, I, p. 19, prop. 4).

For \(x \in A\) and \(t \in I\), we have \(\varphi(x) = 0\), therefore \(t + \varphi(x) \leqslant 1\).

\[
\rho (j (\alpha_ {i} (x, t))) = \rho (x, t) = \alpha_ {i} (\sigma (x, 1), t) = \alpha_ {i} (x, t).
\]

For \(x\in \mathrm{X}\) , we also have

\[
\rho (j (\beta_ {i} (x))) = \rho (x, 1) = \beta_ {i} (\sigma (x, 1)) = \beta_ {i} (x).
\]

As the images of the maps \(\alpha_{i}\) and \(\beta_{i}\) cover \(\mathrm{Cyl}(i)\), it follows that the map \(\rho\) is a continuous retraction of the map j. Consequently, the pair (X, A) has the homotopy extension property (III, p. 241, prop. 7), which concludes the proof of the theorem.

Lemma 4. --- Let X and Y be topological spaces, let \(\varphi\colon\mathrm{X}\to\mathbf{R}_{+}\) be a continuous map, and set \(\mathrm{A}=\overline{\varphi}^{1}(0)\). Let \(\sigma\colon\mathrm{X}\times\mathbf{I}\to\mathrm{Y}\) be a homotopy that is fixed on A. The map \(\sigma'\colon\mathrm{X}\times\mathbf{I}\to\mathrm{Y}\) which maps \((x,s)\) to \(\sigma(x,s/\varphi(x))\) if \(s<\varphi(x)\) and to \(\sigma(x,1)\) if \(s\geqslant\varphi(x)\) is continuous.

The map \(\sigma'\) is continuous at every point of the closed subset \(\overline{\varphi}^{1}([1, +\infty[)\times\mathbf{I}\); it therefore suffices to show that its restriction to \(\overline{\varphi}^{1}(\mathbf{I})\times\mathbf{I}\) is continuous. We may therefore assume that \(\varphi(\mathrm{X})\subset\mathbf{I}\). Let C and C' be the subspaces of X × I consisting of the pairs \((x,s)\) such that \(s\leqslant\varphi(x)\) and \(s\geqslant\varphi(x)\), respectively. They are closed and cover X × I. The map \(\sigma'\) is continuous on C'; let us prove that it is continuous on C.

Let \(\alpha\colon\mathbf{X}\times\mathbf{I}\to\mathbf{X}\times\mathbf{I}\) be the continuous map given by \(\alpha(x,t)=(x,t\varphi(x))\). Its image is equal to C and \(\sigma'\circ\alpha\) is the continuous map \(\sigma\). Let us prove that \(\alpha\) is a proper map. Indeed, consider an ultrafilter \(\mathfrak{U}\) on \(\mathbf{X}\times\mathbf{I}\) and a point \((x,t)\in\mathbf{X}\times\mathbf{I}\) which is adherent to the ultrafilter base \(\alpha(\mathfrak{U})\). Since \(\mathrm{pr}_{1}\circ\alpha=\mathrm{pr}_{1}\), the ultrafilter base \(\mathrm{pr}_{1}(\mathfrak{U})\) on X converges to \(x\). Since I is compact, there exists a point \(s\in\mathbf{I}\) such that the ultrafilter base \(\mathrm{pr}_{2}(\mathfrak{U})\) converges to \(s\). Then \(\mathfrak{U}\) converges to \((x,s)\). Since \(\alpha\) is continuous, the ultrafilter base \(\alpha(\mathfrak{U})\) converges to \((x,s\varphi(x))\). Since I is separated, we have \(s\varphi(x)=t\), whence \(\alpha(x,s)=(x,t)\), so that \(\alpha\) is proper (TG, I, p. 75, th. 1). It then follows from I, p. 18, example 2 and prop. 9 that \(\sigma'\mid\mathrm{C}\) is continuous.

COROLLARY 1. --- Let X be a normal topological space and let A be a closed subspace of X. Suppose there exists a neighborhood V of A in X and a contraction of V onto A, as well as a continuous map \(f: \mathrm{X} \to \mathbf{R}\) such that \(f^{-1}(0) = \mathrm{A}\). Then the pair (X,A) has the homotopy extension property.

Let \(\rho\colon V \times I \to V\) be a contraction of \(V\) onto \(A\). By definition of a normal space (TG, IX, p. 41, definition 1), there exists a continuous map \(g\colon X \to I\) that takes the value 0 on \(A\) and 1 at every point of \(X - V\). Let \(\varphi\colon X \to R\) be the map given by \(\varphi(x) = |f(x)| + g(x)\) for \(x \in X\); it is continuous. We have \(\overline{\varphi}^{1}(0) = A\) and \(\overline{\varphi}^{1}(I) \subset V\). Let \(\sigma\) be the map from \(\overline{\varphi}^{1}(I) \times I\) to \(X\) given by \(\sigma(x,t) = \rho(x,1-t)\) for \(x \in \overline{\varphi}^{1}(I)\) and \(t \in I\). For \(x \in \overline{\varphi}^{1}(I)\), we have \(\sigma(x,1) = \rho(x,0) = x\) and \(\sigma(x,0) = \rho(x,1) \in A\).

The maps \(\varphi\) and \(\sigma\) satisfy the conditions of assertion (iii) of Theorem 1 in III, p. 243; consequently, the pair (X,A) has the homotopy extension property.

Example 1. --- Take as the space X the ball \(\mathbf{B}_{n}\) and for subspace A the sphere \(\mathbf{S}_{n-1}\). If \(V = X - \{0\}\), there exists a strong contraction of V onto \(\mathbf{S}_{n-1}\) (III, p. 237, example). Consequently, the pair \((\mathbf{B}_{n}, \mathbf{S}_{n-1})\) has the homotopy extension property.

COROLLARY 2. --- Let X and Y be topological spaces, let A be a closed subspace of X, and let B be a closed subspace of Y. If the pairs (X,A) and (Y,B) have the homotopy extension property, then so does the pair (X × Y, (X × B) ∪ (A × Y)).

Let \(\varphi\colon X\to\mathbf{R}_{+}\) and \(\sigma\colon\overline{\varphi}^{1}(\mathbf{I})\times\mathbf{I}\to X\), and let \(\varphi^{\prime}\colon Y\to\mathbf{R}_{+}\) and \(\sigma^{\prime}\colon\overline{\psi}^{1}(\mathbf{I})\times\mathbf{I}\to Y\), be maps satisfying the conditions of assertion (iv) of Theorem 1 for the pair (X,A) and, respectively, for the pair (Y,B). Let \(\psi\colon X\times Y\to\mathbf{R}_{+}\) be the map given by \((x,y)\mapsto\inf(\varphi(x),\varphi^{\prime}(x))\); it is continuous; we also have \(\psi(x,y)=0\) for all \((x,y)\in X\times Y\) such that \(x\in A\) or \(y\in B\). For \((x,y,t)\in X\times Y\times I\) such that \(t+\psi(x)\geqslant1\), we have \(t+\varphi(x)\geqslant1\) and \(t+\varphi^{\prime}(x)\geqslant1\). One thus defines a continuous map

\[
\tau \colon \{(x, y, t) \in \mathrm{X} \times \mathrm{Y} \times \mathbf {I} \mid t + \psi (x) \geqslant 1 \} \rightarrow \mathrm{X} \times \mathrm{Y}
\]

by setting \(\tau(x,y,t) = (\sigma(x,t),\sigma'(y,t))\). For all \((x,y)\in X\times Y\), we have \(\tau(x,y,1) = (\sigma(x,1),\sigma'(y,1)) = (x,y)\). If, moreover, \(\psi(x,y)\leqslant 1\), then \(\varphi(x)\leqslant 1\) or \(\varphi'(y)\leqslant 1\), so that \(\sigma(x,1)\in A\) or \(\sigma'(y,1)\in B\); this implies that \(\tau(x,y,1 - \psi(x))\) belongs to \((A\times Y)\cup (X\times B)\). This verifies assertion (iv) of Theorem 1, whence the corollary.

Example 2. --- *Here are other cases where the pair (X, A) possesses the homotopy extension property.

\begin{enumerate}
\def\labelenumi{(\roman{enumi})}
\item
  The space X is a paracompact differentiable manifold of class C \(^{1}\) and A a closed submanifold of X. In view of Corollary \(^{1}\), this follows from the fact that X is perfectly normal (IX, p. 103, exerc. 11) and that A has a tubular neighborhood in X;
\item
  The space X is a cell complex and A a subcomplex relative to a given cell decomposition of X.*
\end{enumerate}
